\documentclass[11pt]{article}
\usepackage[margin=1in]{geometry}
\usepackage{amsmath,amssymb,amsthm,mathtools,bm}
\usepackage{booktabs}
\usepackage[colorlinks=true,linkcolor=blue,citecolor=blue]{hyperref}

\newtheorem{theorem}{Theorem}
\newtheorem{lemma}{Lemma}
\newtheorem{proposition}{Proposition}
\newtheorem{corollary}{Corollary}
\theoremstyle{definition}
\newtheorem{definition}{Definition}
\newtheorem{assumption}{Assumption}
\theoremstyle{remark}
\newtheorem{remark}{Remark}

% macros of multistep_rollout_kernel.tex
\newcommand{\uu}{\mathbf{u}}
\newcommand{\ee}{\mathbf{e}}
\newcommand{\G}{\mathbf{G}}
\newcommand{\Gt}{\tilde{\mathbf{G}}}
\newcommand{\I}{\mathbf{I}}
\newcommand{\Hop}{\mathbf{H}}
\newcommand{\eps}{\bm{\epsilon}}
\newcommand{\Rgh}{\mathcal{R}}
\newcommand{\dt}{\Delta t}
\newcommand{\LL}{\mathcal{L}}
\newcommand{\KK}{\mathcal{K}}
% this note
\newcommand{\Pp}{\mathbf P}
\newcommand{\Ph}{\hat{\mathbf P}}
\newcommand{\LLh}{\hat{\mathcal{L}}}
\newcommand{\Phih}{\hat\Phi}
\newcommand{\Psih}{\hat\Psi}
\newcommand{\etab}{\bm{\eta}}
\newcommand{\Nn}{\mathcal N_\theta}

\title{Multistep training of learned steppers:\\ an expansion to all orders about the true trajectory}
\author{Draft (notation of the rollout-kernel note)}
\date{}

\begin{document}
\maketitle

\begin{abstract}
Multistep (unrolled) training of a learned one-step map is analysed by expanding the rollout about the true trajectory and keeping every order.
The rollout error, the rollout features and the gradient of the multistep loss are written as explicit series in the one-step errors, with coefficients that are derivatives of the learned map along the true trajectory.
For any update rule, the exact error at any set of horizons changes at a rate that is the inner product of its gradient with the parameter update; under gradient flow on a window loss this is an inner product of two gradients, and every order of it is explicit.
If the network is analytic on complex balls around the true states and the first-order errors are small enough relative to the radius (an admissibility condition), the series converge, every order and the remainder after any order are bounded, and, given certified analyticity constants (not yet computed), the sign of the rate can be certified.
Separately, without any expansion: training on a window does not change the method, so the local-complexity lower bound of the roughness theory applies during multistep training, and because the rollout error controls the injected error, it also bounds the window-averaged root-mean-square rollout error from below.
All results are at one training time.
\end{abstract}

\section{What this note is for}\label{sec:purpose}

The roughness theory considers a learned one-step map written as a method: a known part plus a network slot.
Its central result bounds the injected (one-step) error from below by the local complexity of the method (the roughness, on the training inputs, of the target the known part leaves for the network) minus the roughness of what the network currently supplies, $c_p\Nn$, and a leak term set by the training kernel, all divided by the norm constant of the seminorm; the bound holds at every training time and for any update rule, and is informative only when what the network supplies is smoother than its target.
The theory also describes what a single training step can reach: exactly as a bound on the error after the step actually taken, or for every step of a given length through a maximum over a parameter ball, and otherwise to first order in the step length.

Multistep training penalizes rollout errors rather than one-step errors.
This note asks what that changes, and keeps every order of the one-step error throughout: no statement below rests on a truncation of the expansion in the one-step errors, which appears only as a partial sum of a series whose remainder is bounded.
The few statements that are first order in the length of a parameter step say so.
The approach is the error expansion of the eigenanalysis paper, carried to all orders, together with the factorization of the rollout features of the rollout-kernel note.
All quantities are defined from the data (the true trajectories) and the network; several constants used in the bounds (radii of analyticity and sizes on complex balls) are only estimated, as stated where they appear.

\section{The framework in brief}\label{sec:framework}

\paragraph{Setting.}
A learned one-step map $\Gt(\uu;\theta)$ (Euler, $\uu+\dt\,\Hop$, or direct) is trained on a window loss $L_w=\sum_q w_q\tfrac12\|\ee(t_q)\|^2$ over $p$ steps.
Everything is expressed along the \emph{true} trajectory: the one-step errors $\eps_i$, the Jacobians $J_i$ and propagators $\Pp(t_q,t_i)$ of the learned map, and its parameter derivatives $\Phi_i$.
These do not depend on the rollout error, so expansions in the one-step errors have fixed coefficients.

\paragraph{Proved.}
\begin{enumerate}
\item \emph{Expansions to all orders} (Sections~\ref{sec:error}--\ref{sec:grad}). The rollout error, the rollout features and the gradient of $L_w$ are series in $\eps$ with every term explicit. They converge to the exact quantities under the analyticity conditions of Section~\ref{sec:analytic}; without them they are formal.
\item \emph{Exact rates} (Section~\ref{sec:rates}). For any update rule, the exact error at any horizons $v$ changes at the rate $\dot V_v=\langle g_v,\hat g\rangle$; under gradient flow on $L_w$ this is $-\langle g_v,g_w\rangle$. Every order of this rate is explicit.
\item \emph{Convergence, remainder and certified signs} (Sections~\ref{sec:analytic}--\ref{sec:bounds}). If the network is analytic on complex balls of radius $r$ around the true states and an admissibility condition holds, the series converge, every order decays geometrically, the remainder after any order is bounded explicitly, and the sign of the rate can be certified, given certified values of the analyticity constants (not yet computed).
\item \emph{Bounds on the higher-order part} (Section~\ref{sec:bounds}). The exact difference between the error and its first term, and between the gradient and its first term, are bounded to all orders; under uniform constants this gives a window of controlled nonlinearity. That window controls the size of the higher orders; it is not a window in which training is guaranteed to help.
\item \emph{Local-complexity floor} (Section~\ref{sec:floor}). The roughness theory's lower bound on the injected error holds at every time during multistep training, and $\|\eps\|\le(1+\ell)\|\ee\|$ (without expansion; $\ell$ the supremum of $\|\nabla\Gt\|$ on the segments joining true and predicted states) carries it to the window-averaged root-mean-square rollout error. The bound used is the supremum over all roughness functions, the residual $\|E_\gamma\|$ certified by its Krylov values; a fixed seminorm gives a weaker bound. For a direct model, the part of the identity map that a step cannot build bounds the rollout error after the step from below, up to a term of order $\dt$.
\item \emph{Operator errors} (Section~\ref{sec:operators}). The rollout error is exactly the separation of true and predicted states by the true flow plus the model error at the predicted state. Each state derivative of the learned map splits into a derivative of the true flow map and an error of the learned operator; along the true trajectory, the learned propagators differ from the true tangent flow by an amount bounded through the Jacobian errors alone, given the true-system growth constants. Each operator error (values, Jacobian, higher derivatives) obeys the roughness theory with its own features, with supremum and Krylov certificates; one-step training reaches it only through a cross-kernel. With all coefficients held fixed, the total norm of small injections cannot decide whether a window helps the error at other horizons; under an admissibility condition, the expected rate (the first-order part of the expected one-step decrease that defines the window for one training step) is a series of joint moments of injections, operator errors and features.
\item \emph{One training step, to any order in the step} (Section~\ref{sec:stepseries}). Along a direction in parameter space, the rollout and the window error are power series in the step length, with sensitivities of every order given by an exact recursion along the rollout and a bounded remainder after any order inside a guaranteed disc of radius $R_p$. The roughness paper's single-step value is the minimum over directions of the model with the error linearized in the step; the series adds the rest along a given direction, so inside its radius the step is described exactly to any order, and the radius, not the nonlinearity, limits the step lengths that any expansion about the current parameters can describe. This certifies when one step of a given length reduces the window's own error, and so gives the largest window for which such a step is certified (a limit of certification, with no converse). For Euler models the guaranteed radius shrinks at most exponentially in the physical length of the window, independently of $\dt$ for given network constants.
\item \emph{The two radii} (Section~\ref{sec:tworadii}). For SiLU MLPs the radius of convergence of the step series is bounded from below by keeping every hidden pre-activation away from the poles of the activation, and from above, via Rouché's theorem and under hypotheses on complex discs that are not yet verified for trained networks, by a last-layer pre-activation that reaches a pole (in a conjugate pair) near its linear prediction; the upper bound applies to the rollout state one step later, and to the window error when the window ends there. To first order both bounds are the distance of a pre-activation from the nearest pole divided by how fast the step moves it, so the radius in parameter space is, to first order, at least the radius in state space divided by how far a unit step moves the predicted state.
\end{enumerate}

\paragraph{Properties of the network, not of the system.}
The growth of the propagators, the radius $r$, the sizes on complex balls and the Lipschitz constant $\ell$ are properties of the trained network.
They approximate properties of the system only to the extent that the network approximates the flow map and its derivatives; Section~\ref{sec:operators} makes this exact for the propagators along the true trajectory, separating the true tangent flow from the Jacobian errors, and bounds $D_i$, $D'_i$ by their true and error parts.

\paragraph{Not proved.}
How far a rollout window can usefully be extended is not determined by these results: they certify, at one training time, that training on a window decreases the error at chosen horizons whenever the certified sign is positive, and they cannot prove the converse.
The growth of the Gauss--Newton part of the loss curvature with the window is proved only as an upper bound (rollout-kernel note); its observed exponential growth is measured.
The concentration of the gradient on a few starts is measured.

\paragraph{Gaps.}
\begin{enumerate}
\item \emph{What training reaches.} All results are about one training time; nothing integrates over a run.
\item \emph{Minibatches.} A minibatch step inserts a random sampling matrix between the kernel and the error, $\dot\ee=-\Psih\Psih^\top SW\ee$ in continuous time; no statement is made about the random product. For a discrete step $\Delta\theta$ the exact change is $\langle\bar g_v,\Pi^{-1/2}\Delta\theta\rangle$ with $\bar g_v$ averaged along the step; Lemma~\ref{lem:rates} is its continuous-time form.
\item \emph{Adaptive optimizers.} The same holds with Adam's actual step; there is no closed form linking Adam's direction to the current gradient.
\item \emph{Beyond the radius.} Where a start's error exceeds the radius of analyticity, the series need not converge; the exact identities (Lemma~\ref{lem:recursion}, Lemma~\ref{lem:rates}, Section~\ref{sec:floor}) still hold.
\item \emph{A floor at a fixed horizon.} No local-complexity floor is known for the error at a fixed horizon in terms of the roughness of the flow map over that horizon.
\end{enumerate}

\section{Setting and notation}\label{sec:setting}

The learned one-step map is $\Gt(\uu;\theta)$ with a smooth activation; for forward Euler $\Gt(\uu)=\uu+\dt\,\Hop(\uu;\theta)$, for a direct model $\Gt$ is the network itself.
The data are noise-free samples of true trajectories at spacing $\dt$.
Training rollouts start from $n$ true states and have $p$ steps.
For each start (index $k$, suppressed where it is not needed):
\begin{itemize}
\item true states $\uu_T(t_i)$, predicted states $\uu_p(t_0)=\uu_T(t_0)$, $\uu_p(t_{i+1})=\Gt(\uu_p(t_i))$;
\item rollout errors $\ee(t_i)=\uu_T(t_i)-\uu_p(t_i)$, $\ee(t_0)=0$;
\item one-step errors at the true states, $\eps_i=\uu_T(t_{i+1})-\Gt(\uu_T(t_i))$, $i=0,\dots,p-1$; for Euler models $\eps_i=\dt\,\etab_i$;
\item Jacobians of the learned map along the \emph{true} trajectory, $J_i=\nabla\Gt(\uu_T(t_i))$, and propagators $\Pp(t_q,t_i)=J_{q-1}\cdots J_{i+1}$, $\Pp(t_{i+1},t_i)=\I$;
\item the same along the \emph{predicted} trajectory, $\hat J_i=\nabla\Gt(\uu_p(t_i))$, $\Ph(t_q,t_i)=\hat J_{q-1}\cdots\hat J_{i+1}$;
\item parameter derivatives $\Phi_i=\nabla_\theta\Gt(\uu_T(t_i))\Pi^{1/2}$ at the true states and $\Phih_i=\nabla_\theta\Gt(\uu_p(t_i))\Pi^{1/2}$ at the predicted states, with $\Pi$ a fixed diagonal preconditioner with positive entries.
\end{itemize}
Stacking over starts and steps, $\LL$ maps injections $z=(z_i)_{i<p}$ to $(\LL z)_q=\sum_{i<q}\Pp(t_q,t_i)z_i$; $\LLh$ is the same with $\Ph$.
The multistep loss is $L_w=\tfrac12\,\ee^\top W\ee$, where $W$ contains weights $w_q\ge0$ and the average $1/n$ over starts ($w_q=1/p$ for the time-averaged loss).
The rollout features are $\Psih=\partial_\theta\uu_p\,\Pi^{1/2}$, and by the features lemma of the rollout-kernel note $\Psih=\LLh\Phih$ exactly; along the true trajectory write $\Psi=\LL\Phi$.
The rescaled parameter velocity is $\hat g=\Pi^{-1/2}\dot\theta$; for a loss $X$ write $g_X=\Pi^{1/2}\nabla_\theta X$, so gradient flow on $L_w$ is $\hat g=-g_w$ with $g_w=-\Psih^\top W\ee$.
The rollout-kernel note puts the factor $1/n$ in its kernel; here it is in $W$.

\begin{remark}[Hats]
In the rollout-kernel note $\Phi$ and $\LL$ are evaluated along the predicted trajectory; here those carry a hat, and unhatted quantities are evaluated along the true trajectory, whose coefficients do not depend on the error.
\end{remark}

\section{The rollout error to all orders}\label{sec:error}

\begin{lemma}[Error recursion]\label{lem:recursion}
Exactly, $\ee(t_{i+1})=J_i\ee(t_i)+\eps_i+N_i(\ee(t_i))$ with $N_i(x)=\Gt(\uu_T(t_i))-\Gt(\uu_T(t_i)-x)-J_ix$.
If $\Gt$ is analytic and $|\ee(t_i)|$ is within the radius of convergence of its Taylor series at $\uu_T(t_i)$ in the direction of $\ee(t_i)$, then its Taylor series converges to $N_i$:
\begin{equation}\label{eq:recursion}
N_i(\ee(t_i))=\sum_{k\ge2}\frac{(-1)^{k+1}}{k!}\,\nabla^k\Gt\big(\uu_T(t_i)\big)\big[\ee(t_i)^{\otimes k}\big].
\end{equation}
\end{lemma}

\begin{proof}
$\ee(t_{i+1})=\uu_T(t_{i+1})-\Gt(\uu_p(t_i))=\eps_i+\Gt(\uu_T(t_i))-\Gt(\uu_T(t_i)-\ee(t_i))$; expand the last term about $\uu_T(t_i)$.
\end{proof}

\begin{definition}[Orders of the error]\label{def:orders}
$\ee^{(1)}=\LL\eps$ and, for $m\ge2$,
\begin{equation}\label{eq:orders}
\ee^{(m)}(t_q)=\sum_{i<q}\Pp(t_q,t_i)\sum_{k=2}^{m}\frac{(-1)^{k+1}}{k!}\sum_{\substack{m_1+\dots+m_k=m\\ m_l\ge1}}\nabla^k\Gt\big(\uu_T(t_i)\big)\big[\ee^{(m_1)}(t_i),\dots,\ee^{(m_k)}(t_i)\big].
\end{equation}
\end{definition}

$\ee^{(m)}$ is a homogeneous polynomial of degree $m$ in $\eps$ whose coefficients are derivatives of $\Gt$ along the true trajectory; for example $\ee^{(2)}(t_q)=-\tfrac12\sum_{i<q}\Pp(t_q,t_i)\nabla^2\Gt(\uu_T(t_i))[\ee^{(1)}(t_i),\ee^{(1)}(t_i)]$.
Inserting $\ee=\sum_m\ee^{(m)}$ into Lemma~\ref{lem:recursion} and collecting terms of equal degree reproduces \eqref{eq:orders}, so $\sum_m\ee^{(m)}$ is the formal power series of the error in the one-step errors.
Theorem~\ref{thm:convergence} proves that it converges to $\ee$ under explicit conditions; until then the identity $\ee=\sum_m\ee^{(m)}$ is formal.

\begin{remark}[Radius of convergence]\label{rem:radius}
A network with SiLU (or tanh) activation is analytic on the real line but has complex singularities.
For the width-256 Lorenz-63 models the one-variable radius of $\zeta\mapsto\Gt(\uu+\zeta\mathbf n)$ was estimated along real directions $\mathbf n$ at held-out states (Taylor coefficients from samples on complex circles): median $9$--$13$ state units, minimum $4.3$ over the sample.
This is an upper bound for the radius $r$ of Assumption~\ref{ass:analytic} below, which requires holomorphy on a full complex ball, and a sampled minimum is not a certified value.
\end{remark}

\section{The training gradient to all orders}\label{sec:grad}

\paragraph{Exact form.}
$\partial_\theta\ee=-\Psih\,\Pi^{-1/2}$, so $g_w=-\Psih^\top W\ee$ exactly, with $\Psih=\LLh\Phih$.

\paragraph{Expansion of the predicted-trajectory quantities.}
Since $\Phih_i=\nabla_\theta\Gt(\uu_T(t_i)-\ee(t_i))\Pi^{1/2}$ and $\hat J_i=\nabla\Gt(\uu_T(t_i)-\ee(t_i))$, expanding about the true states and inserting the orders of the error gives formal series $\Phih_i=\sum_{j\ge0}\Phi^{(j)}_i$, $\hat J_i=\sum_{j\ge0}J^{(j)}_i$ with $\Phi^{(0)}=\Phi$, $J^{(0)}=J$ and, for $j\ge1$,
\begin{align}
\Phi^{(j)}_i&=\sum_{k=1}^{j}\frac{(-1)^k}{k!}\sum_{\substack{m_1+\dots+m_k=j\\ m_l\ge1}}\nabla_\uu^k\nabla_\theta\Gt\big(\uu_T(t_i)\big)\big[\ee^{(m_1)}(t_i),\dots,\ee^{(m_k)}(t_i)\big]\,\Pi^{1/2},\label{eq:Phij}\\
J^{(j)}_i&=\sum_{k=1}^{j}\frac{(-1)^k}{k!}\sum_{\substack{m_1+\dots+m_k=j\\ m_l\ge1}}\nabla^{k+1}\Gt\big(\uu_T(t_i)\big)\big[\ee^{(m_1)}(t_i),\dots,\ee^{(m_k)}(t_i),\,\cdot\,\big].\label{eq:Jj}
\end{align}
Products of Jacobians expand by distributing the orders over the factors,
\[
\Pp^{(j)}(t_q,t_i)=\sum_{\substack{j_{i+1}+\dots+j_{q-1}=j\\ j_l\ge0}}J^{(j_{q-1})}_{q-1}\cdots J^{(j_{i+1})}_{i+1},\qquad (\LL^{(j)}z)_q=\sum_{i<q}\Pp^{(j)}(t_q,t_i)\,z_i ,
\]
with the empty product for $q=i+1$ giving $\Pp^{(0)}=\I$ and $\Pp^{(j)}=0$ for $j\ge1$,
and, since $\Psih=\LLh\Phih$,
\begin{equation}\label{eq:Psij}
\Psih=\sum_{j\ge0}\Psi^{(j)},\qquad \Psi^{(j)}=\sum_{a+b=j}\LL^{(a)}\Phi^{(b)},\qquad \Psi^{(0)}=\Psi .
\end{equation}

\paragraph{The gradient by order.}
Collecting terms of total degree in $\eps$ in $g_w=-\Psih^\top W\ee$,
\begin{equation}\label{eq:gradorders}
g_w=\sum_{m\ge1}\nabla^{(m)}_w,\qquad \nabla^{(m)}_w=-\sum_{j=0}^{m-1}\big(\Psi^{(j)}\big)^\top W\,\ee^{(m-j)},
\end{equation}
again formal until Theorem~\ref{thm:convergence}.
The term of degree one is $\nabla^{(1)}_w=-\Phi^\top\LL^\top W\LL\eps$; the term of degree two is $-\Psi^\top W\ee^{(2)}-(\LL\Phi^{(1)}+\LL^{(1)}\Phi)^\top W\ee^{(1)}$, with $\Phi^{(1)}_i=-\nabla_\uu\nabla_\theta\Gt(\uu_T(t_i))[\ee^{(1)}(t_i)]\Pi^{1/2}$.

\paragraph{Where the change of the learned Jacobian enters.}
The mixed derivatives $\nabla_\uu^k\nabla_\theta\Gt$ (how a parameter change alters the learned map away from the true states, and hence its Jacobian) appear only through $\Phi^{(j)}$, $j\ge1$, and therefore only in $\nabla^{(m)}_w$ with $m\ge2$.
By symmetry of mixed partial derivatives, $\Phi^{(1)}_i\delta\theta=-\big(\partial_\theta J_i\,\delta\theta\big)\ee^{(1)}(t_i)$ for a parameter change $\delta\theta$ (with $\Pi^{1/2}$ absorbed): these are the terms through which the gradient takes into account how the propagation of errors depends on the parameters.

\section{Exact rates at any horizon}\label{sec:rates}

Let $v=(v_q)$ be nonnegative weights on horizons $t_q$, $q=1,\dots,\bar p$, where the rollout is extended to $\bar p\ge p$ steps so that it covers the evaluation horizons ($w_q=0$ for $q>p$), and
\[
V_v=\tfrac12\,\ee^\top W_v\,\ee=\tfrac1n\sum_k\sum_q v_q\,\tfrac12|\ee_k(t_q)|^2 ,
\]
the exact error at those horizons; a single horizon $t_{q'}$ is $v_q=\delta_{qq'}$.
Then $g_v=-\Psih^\top W_v\ee$.

\begin{lemma}[Exact rates]\label{lem:rates}
For any update rule with rescaled parameter velocity $\hat g$, $\dot V_v=\langle g_v,\hat g\rangle$.
Under gradient flow on $L_w$,
\begin{equation}\label{eq:tradeoff}
\dot V_v=-\langle g_v,g_w\rangle=-\big(W_v\ee\big)^\top\Psih\Psih^\top\big(W\ee\big).
\end{equation}
\end{lemma}

\begin{proof}
$\dot V_v=\nabla_\theta V_v^\top\dot\theta=(\Pi^{1/2}\nabla_\theta V_v)^\top(\Pi^{-1/2}\dot\theta)$.
\end{proof}

No expansion is involved.
Training on window $w$ decreases the error at horizons $v$ at this instant if and only if $\langle g_v,g_w\rangle>0$; for $v=w$ the rate is $-|g_w|^2\le0$, so the content is in horizons other than the window's own loss.

\begin{proposition}[Every order of the rate]\label{prop:rateorders}
As formal series, $\dot V_v=-\sum_{M\ge2}R^{(M)}$ with
\begin{equation}\label{eq:RM}
R^{(M)}=\sum_{\substack{a+j+l+b=M\\ a,b\ge1;\ j,l\ge0}}\big(W_v\ee^{(a)}\big)^\top\Psi^{(j)}\big(\Psi^{(l)}\big)^\top\big(W\ee^{(b)}\big),
\end{equation}
homogeneous of degree $M$ in $\eps$; for instance $R^{(2)}=\langle\nabla^{(1)}_v,\nabla^{(1)}_w\rangle$.
If the hypothesis of Theorem~\ref{thm:convergence} holds for all starts, some $\rho>1$ is admissible (the admissible set is open), and by Theorem~\ref{thm:tradeoffbound} the series converges to $\dot V_v$ and the remainder after any order is bounded.
\end{proposition}

\begin{proof}
Insert the series for $\ee$ and $\Psih$ into \eqref{eq:tradeoff} and collect terms of degree $M$; there are no terms with $a=0$ or $b=0$ because $\ee$ has no degree-zero part.
\end{proof}

\section{Analyticity, convergence and the remainder}\label{sec:analytic}

\begin{assumption}[Analytic extension]\label{ass:analytic}
There is $r>0$ such that for every true state $\uu_i:=\uu_T(t_i)$ of the rollout the maps
\[
g_i(z)=\Gt(\uu_i+z)-\Gt(\uu_i)-z,\qquad h_i(z)=\nabla_\theta\Gt(\uu_i+z)\,\Pi^{1/2}
\]
extend holomorphically to $\{z\in\mathbb C^d:|z|<r\}$ ($|\cdot|$ the Hermitian norm), with
\[
D_i=\sup_{|z|<r}|g_i(z)|,\qquad D_i'=\sup_{|z|<r}\|\nabla g_i(z)-\nabla g_i(0)\|,\qquad F_i=\sup_{|z|<r}\|h_i(z)-h_i(0)\|
\]
finite.
For forward Euler $g_i(z)=\dt\,(\Hop(\uu_i+z)-\Hop(\uu_i))$, so $D_i=\dt\,D_{H,i}$ and $D_i'=\dt\,D'_{H,i}$.
\end{assumption}

The constants blow up as $r$ approaches the true radius of holomorphy, so $r$ must be taken strictly inside it.
Operator norms of real matrices are the same over $\mathbb R$ and $\mathbb C$.
The Taylor coefficients of order $k\ge2$ of $\Gt$ at $\uu_i$ are those of $g_i$, since the two differ by an affine map.
Write $\varphi(s)=s^2/(1-s)$ and $\psi(s)=s/(1-s)$ for $0\le s<1$; both are increasing.

\begin{lemma}[Cauchy estimates]\label{lem:cauchy}
Under Assumption~\ref{ass:analytic}, for $x\in\mathbb C^d$ with $|x|<r$:
\begin{align}
\Big|\tfrac1{k!}\nabla^k\Gt(\uu_i)[x^{\otimes k}]\Big|&\le D_i\,(|x|/r)^k\quad(k\ge2),\qquad
|N_i(x)|\le D_i\,\varphi(|x|/r),\label{eq:cauchyN}\\
\big\|\nabla\Gt(\uu_i-x)-J_i\big\|&\le D_i'\,\psi(|x|/r),\qquad
\big\|\nabla_\theta\Gt(\uu_i-x)\Pi^{1/2}-\Phi_i\big\|\le F_i\,\psi(|x|/r),\label{eq:cauchyJ}
\end{align}
with $N_i$ continued holomorphically, and for $k\ge2$ and $x_1,\dots,x_k\in\mathbb R^d$
\begin{equation}\label{eq:cauchymulti}
\Big|\tfrac1{k!}\nabla^k\Gt(\uu_i)[x_1,\dots,x_k]\Big|\le\frac{D_i}{r^k}\prod_j|x_j| .
\end{equation}
\end{lemma}

\begin{proof}
For a unit vector $v\in\mathbb C^d$, $\zeta\mapsto g_i(\zeta v)$ is holomorphic on $|\zeta|<r$ and bounded by $D_i$; Cauchy's estimate on $|\zeta|=\rho<r$ and $\rho\to r$ give the first bound.
$N_i(x)=-\sum_{k\ge2}\tfrac1{k!}\nabla^kg_i(0)[(-x)^{\otimes k}]$, and summing gives the second.
The bounds \eqref{eq:cauchyJ} are the same argument for the matrix- and operator-valued maps $\nabla g_i(z)-\nabla g_i(0)$ and $h_i(z)-h_i(0)$, whose Taylor series start at order one.
For \eqref{eq:cauchymulti}: for a unit $w\in\mathbb R^d$, $\langle w,\tfrac1{k!}\nabla^k\Gt(\uu_i)[\cdot]\rangle$ is a symmetric $k$-linear form on $\mathbb R^d$, whose norm equals the supremum of its diagonal on the unit sphere (Banach's theorem on symmetric multilinear forms on Hilbert spaces); the first bound controls the diagonal; take the supremum over $w$.
\end{proof}

\begin{theorem}[The higher-order part of the error]\label{thm:allorders}
Under Assumption~\ref{ass:analytic}, for each start define recursively $a_0=0$ and, while the arguments of $\varphi$ are below one,
\[
a_q=\sum_{i<q}\|\Pp(t_q,t_i)\|\,D_i\,\varphi\Big(\frac{|\ee^{(1)}(t_i)|+a_i}{r}\Big).
\]
If $|\ee^{(1)}(t_i)|+a_i<r$ for all $i<q$, then $|\ee(t_i)-\ee^{(1)}(t_i)|\le a_i$ for all $i\le q$.
\end{theorem}

\begin{proof}
By Lemma~\ref{lem:recursion}, $\bm\delta=\ee-\ee^{(1)}$ satisfies $\bm\delta(t_{i+1})=J_i\bm\delta(t_i)+N_i(\ee(t_i))$, $\bm\delta(t_0)=0$, so $\bm\delta(t_q)=\sum_{i<q}\Pp(t_q,t_i)N_i(\ee(t_i))$.
If $|\bm\delta(t_i)|\le a_i$ for $i<q$, then $|\ee(t_i)|<r$, \eqref{eq:cauchyN} applies, and summing gives $|\bm\delta(t_q)|\le a_q$; induction on $q$.
\end{proof}

This bounds the exact $\ee-\ee^{(1)}$; it does not use the series.

\begin{theorem}[Convergence of the error series]\label{thm:convergence}
Under Assumption~\ref{ass:analytic} and the hypothesis of Theorem~\ref{thm:allorders} for all $q\le\bar p$, the series $\sum_m\ee^{(m)}(t_q)$ converges absolutely to $\ee(t_q)$, with $\sum_{m\ge2}|\ee^{(m)}(t_q)|\le a_q$.
\end{theorem}

\begin{proof}
Let $B^{(1)}_i=|\ee^{(1)}(t_i)|$ and, for $m\ge2$, $B^{(m)}_q=\sum_{i<q}\|\Pp(t_q,t_i)\|D_i\sum_{k=2}^mr^{-k}\sum_{m_1+\dots+m_k=m,\ m_l\ge1}\prod_lB^{(m_l)}_i$.
By \eqref{eq:orders}, \eqref{eq:cauchymulti} and induction on $m$, $|\ee^{(m)}(t_q)|\le B^{(m)}_q$.
Let $A^{\le M}_q=\sum_{m=2}^MB^{(m)}_q$. Every product in $B^{(m)}_q$ with $m\le M$ involves only orders below $m$, and $\sum_{k\ge2}(\beta/r)^k=\varphi(\beta/r)$ for $\beta<r$, so $A^{\le M}_q\le\sum_{i<q}\|\Pp(t_q,t_i)\|D_i\varphi\big((B^{(1)}_i+A^{\le M-1}_i)/r\big)$; induction on $M$ and $q$ gives $A^{\le M}_q\le a_q$ for all $M$.
Hence $\sum_m|\ee^{(m)}(t_q)|\le|\ee^{(1)}(t_q)|+a_q$, which is below $r$ for $q<\bar p$.
With this absolute convergence, \eqref{eq:cauchymulti} lets the series $N_i(\sum_m\ee^{(m)})=\sum_k\frac{(-1)^{k+1}}{k!}\nabla^k\Gt[(\sum_m\ee^{(m)})^{\otimes k}]$ be regrouped by degree, so $\sum_m\ee^{(m)}$ satisfies the exact recursion of Lemma~\ref{lem:recursion} with the same initial value; the recursion has a unique solution, which is $\ee$.
\end{proof}

\paragraph{Scaling the one-step errors.}
For the features and the rate, scale the one-step errors at a fixed training time: for $s\in\mathbb C$ let $\ee(s)$ solve $\ee(s)(t_{i+1})=s\eps_i+\Gt(\uu_T(t_i))-\Gt(\uu_T(t_i)-\ee(s)(t_i))$, $\ee(s)(t_0)=0$, with $\Gt$ continued holomorphically, and define $\Psih(s)=\LLh(s)\Phih(s)$ with $\hat J_i(s)=\nabla\Gt(\uu_T(t_i)-\ee(s)(t_i))$ and $\Phih_i(s)=\nabla_\theta\Gt(\uu_T(t_i)-\ee(s)(t_i))\Pi^{1/2}$.
At $s=1$ these are the actual error and rollout features.
For $\rho>0$ let $a_{k,q}(\rho)$ be the sequence of Theorem~\ref{thm:allorders} computed with $|\ee^{(1)}_k|$ replaced by $\rho|\ee^{(1)}_k|$, and call $\rho$ \emph{admissible} if $\rho|\ee^{(1)}_k(t_i)|+a_{k,i}(\rho)<r$ for all starts $k$ and steps $i<\bar p$; the set of admissible $\rho$ is open.
Write $\bar e_{k,q}=\rho|\ee^{(1)}_k(t_q)|+a_{k,q}(\rho)$.

\begin{theorem}[Convergence and remainder of the rate]\label{thm:tradeoffbound}
Under Assumption~\ref{ass:analytic}, let $\rho>1$ be admissible. Then:
\begin{enumerate}
\item $s\mapsto\ee(s)$, $\Psih(s)$ and $R(s)=(W_v\ee(s))^\top\Psih(s)\Psih(s)^\top(W\ee(s))$ are holomorphic on a disc containing $|s|\le\rho$; their Taylor coefficients at $s=0$ are $\ee^{(m)}$, $\Psi^{(j)}$ and $R^{(M)}$, and their Taylor series converge at $s=1$ to $\ee$, $\Psih$ and $-\dot V_v$.
\item $|\ee^{(m)}_k(t_q)|\le\bar e_{k,q}\,\rho^{-m}$ for every $m\ge1$: every order decays geometrically.
\item With $B(\rho)\ge\sup_{|s|=\rho}|R(s)|$, for every $N\ge2$
\begin{equation}\label{eq:tradeoffremainder}
\Big|\dot V_v+\sum_{M=2}^{N}R^{(M)}\Big|\le\frac{B(\rho)\,\rho^{-N}}{\rho-1}.
\end{equation}
A computable choice is $B(\rho)=\Big(\frac1n\sum_k\sum_qv_q\,\bar e_{k,q}\bar\Psi_{k,q}\Big)\Big(\frac1n\sum_k\sum_qw_q\,\bar e_{k,q}\bar\Psi_{k,q}\Big)$, the product of the two separately averaged factors, with $\bar\Psi_{k,q}=\sum_{i<q}\overline P_{k}(t_q,t_i)\big(\|\Phi_{k,i}\|+F_{k,i}\psi(\bar e_{k,i}/r)\big)$ and $\overline P_k(t_q,t_i)$ any bound on $\|\Ph_k(s)(t_q,t_i)\|$ over $|s|\le\rho$, for example the Grönwall bound \eqref{eq:gronwall} with $s_j$ replaced by $\bar e_{k,j}/r$ (the same telescoping holds for complex $s$, with $\|\hat J_j(s)-J_j\|\le D'_j\psi(\bar e_{k,j}/r)$).
\item In particular $V_v$ decreases at this instant if $\sum_{M=2}^NR^{(M)}>B(\rho)\rho^{-N}/(\rho-1)$ for some $N$. Since also $|R^{(M)}|\le B(\rho)\rho^{-M}$, the criterion with a given $N$ can hold only if $\rho^{N-1}>2$; with $N=2$ this needs $\rho>2$.
\end{enumerate}
\end{theorem}

\begin{proof}
The Cauchy estimates hold for complex arguments, so the proof of Theorem~\ref{thm:allorders} applies to the injections $s\eps$ for $|s|\le\rho$, giving $|\ee_k(s)(t_q)|\le\bar e_{k,q}<r$; the recursion is then a finite composition of holomorphic maps, so $\ee(s)$, $\hat J_i(s)=\I+\nabla g_i(-\ee(s)(t_i))$, $\Phih_i(s)=h_i(-\ee(s)(t_i))$, $\Psih(s)$ and the bilinear (not Hermitian) form $R(s)$ are holomorphic; admissibility of some $\rho'>\rho$ (openness) gives holomorphy on a disc containing $|s|\le\rho$.
Near $s=0$, $\ee(s)=O(s)$, so collecting powers of $s$ in the recursion and in the products is legitimate and reproduces \eqref{eq:orders}, \eqref{eq:Phij}--\eqref{eq:Psij} and \eqref{eq:RM}; by uniqueness of Taylor coefficients these are the coefficients.
A function holomorphic on a disc of radius above one has a Taylor series converging at $s=1$.
Cauchy's estimate on $|s|=\rho$ gives $|\ee^{(m)}|\le\bar e\rho^{-m}$ and $|R^{(M)}|\le B(\rho)\rho^{-M}$; summing over $M>N$ gives \eqref{eq:tradeoffremainder}.
For $B(\rho)$: $|(W_v\ee)^\top\Psih\Psih^\top(W\ee)|\le|\Psih^\top W_v\ee|\,|\Psih^\top W\ee|$, and each factor is bounded by the triangle inequality over starts and horizons with $\|\Phih_{k,i}(s)\|\le\|\Phi_{k,i}\|+F_{k,i}\psi(\bar e_{k,i}/r)$; the kernel couples starts, so the two factors are averaged separately.
\end{proof}

The same argument applies to the gradient: under the hypotheses of Theorem~\ref{thm:tradeoffbound}, $g_w(s)=-\Psih(s)^\top W\ee(s)$ is holomorphic on a disc of radius above one with Taylor coefficients $\nabla^{(m)}_w$, so the series \eqref{eq:gradorders} converges to $g_w$.

\section{Bounds on the higher-order part}\label{sec:bounds}

\paragraph{The gradient.}
\begin{theorem}[Bound on $g_w-\nabla^{(1)}_w$]\label{thm:gradbound}
Under the hypothesis of Theorem~\ref{thm:allorders} for every start $k$, let $s_{k,i}=(|\ee^{(1)}_k(t_i)|+a_{k,i})/r$ and
\[
\Delta_{k,q}=\sum_{i<q}\Big(\|\Ph_k(t_q,t_i)\|\,F_{k,i}\,\psi(s_{k,i})+\|\Phi_{k,i}\|\sum_{i<j<q}\|\Ph_k(t_q,t_j)\|\,D'_{k,j}\,\psi(s_{k,j})\,\|\Pp_k(t_j,t_i)\|\Big).
\]
Then
\begin{equation}\label{eq:gradbound}
\big|g_w-\nabla^{(1)}_w\big|\le\frac1n\sum_{k}\sum_{q}w_q\Big(\|\Psih_{k,q}\|\,a_{k,q}+\Delta_{k,q}\,|\ee^{(1)}_k(t_q)|\Big).
\end{equation}
\end{theorem}

\begin{proof}
Exactly, $g_w-\nabla^{(1)}_w=-\Psih^\top W\ee+\Psi^\top W\ee^{(1)}=-\Psih^\top W\bm\delta-(\Psih-\Psi)^\top W\ee^{(1)}$ with $\bm\delta=\ee-\ee^{(1)}$.
Per start and horizon, $\Psih_{k,q}-\Psi_{k,q}=\sum_{i<q}\big(\Ph(t_q,t_i)(\Phih_i-\Phi_i)+(\Ph(t_q,t_i)-\Pp(t_q,t_i))\Phi_i\big)$ and the product telescopes exactly, $\Ph(t_q,t_i)-\Pp(t_q,t_i)=\sum_{i<j<q}\Ph(t_q,t_j)(\hat J_j-J_j)\Pp(t_j,t_i)$.
Bound $\|\Phih_i-\Phi_i\|$ and $\|\hat J_j-J_j\|$ by \eqref{eq:cauchyJ} with $|\ee(t_i)|\le s_ir$, $|\bm\delta_k(t_q)|$ by $a_{k,q}$, and use the triangle inequality ($w_q\ge0$).
\end{proof}

\begin{corollary}[Certified sign of the rate]\label{cor:sign}
Let $b_v$, $b_w$ be the right sides of \eqref{eq:gradbound} for the losses $V_v$ and $L_w$. Then $V_v$ decreases at this instant under gradient flow on $L_w$ if
\[
\langle\nabla^{(1)}_v,\nabla^{(1)}_w\rangle>|\nabla^{(1)}_v|\,b_w+b_v\,|\nabla^{(1)}_w|+b_vb_w .
\]
\end{corollary}

\begin{proof}
Write $g_v=\nabla^{(1)}_v+\bm\delta_v$, $g_w=\nabla^{(1)}_w+\bm\delta_w$ with $|\bm\delta_v|\le b_v$, $|\bm\delta_w|\le b_w$, expand $\langle g_v,g_w\rangle$ and use Lemma~\ref{lem:rates}.
\end{proof}

This certificate uses the leading term of the exact rate and a bound on everything else; it needs only the hypothesis of Theorem~\ref{thm:allorders}, not $\rho>2$.
The quantities $\|\Ph\|$ and $\|\Psih\|$ are evaluated along the predicted trajectory.
An a~priori bound follows from the reverse telescoping $\Ph(t_q,t_i)=\Pp(t_q,t_i)+\sum_{i<j<q}\Pp(t_q,t_j)(\hat J_j-J_j)\Ph(t_j,t_i)$ and the discrete Grönwall inequality: if $\|\Pp(t_j,t_i)\|\le Ce^{\Lambda(t_j-t_i)}$ for all $i<j$, then
\begin{equation}\label{eq:gronwall}
\|\Ph(t_q,t_i)\|\le Ce^{\Lambda(t_q-t_i)}\prod_{i<j<q}\big(1+C\,D'_j\psi(s_j)\big),
\end{equation}
and $\|\Psih_{k,q}\|\le\sum_{i<q}\|\Ph_k(t_q,t_i)\|(\|\Phi_{k,i}\|+F_{k,i}\psi(s_{k,i}))$.

\paragraph{A window of controlled nonlinearity.}
\begin{corollary}[Forward Euler, uniform constants]\label{cor:closed}
Suppose that for all starts $\|\Pp(t_j,t_i)\|\le C e^{\Lambda(t_j-t_i)}$ for $i<j$ ($C\ge1$, $\Lambda>0$), $|\etab_i|\le\bar\eta$ and $D_{H,i}\le D_H$, with $t_0=0$ for every start.
Let $\hat E(t)=C\bar\eta\,e^{\Lambda\dt}e^{\Lambda t}/\Lambda$ and $\kappa_r=8CD_H/(r^2\Lambda)$.
If $\hat E(\tau)\le\min(r/4,1/\kappa_r)$, then for every $t_q\le\tau$
\[
|\ee^{(1)}(t_q)|\le\hat E(t_q),\qquad |\ee(t_q)-\ee^{(1)}(t_q)|\le\kappa_r\,\hat E(t_q)^2\le\hat E(t_q).
\]
This holds up to
\[
\tau_{\rm nl}=\frac1\Lambda\ln\frac{\Lambda\,\min(r/4,\,1/\kappa_r)}{C\bar\eta\,e^{\Lambda\dt}} .
\]
\end{corollary}

\begin{proof}
$|\ee^{(1)}(t_q)|\le C\bar\eta\dt\sum_{i<q}e^{\Lambda(t_q-t_i)}\le C\bar\eta e^{\Lambda t_q}\dt/(1-e^{-\Lambda\dt})\le\hat E(t_q)$, using $1-e^{-x}\ge xe^{-x}$.
Show $a_q\le\kappa_r\hat E(t_q)^2\le\hat E(t_q)$ by induction: if $a_i\le\hat E(t_i)$ for $i<q$, then $(|\ee^{(1)}(t_i)|+a_i)/r\le2\hat E(t_i)/r\le1/2$, where $\varphi(s)\le2s^2$, so $a_q\le\frac{8CD_H}{r^2}\dt\sum_{i<q}e^{\Lambda(t_q-t_i)}\hat E(t_i)^2\le\kappa_r\hat E(t_q)^2$, using $\dt\sum_{i<q}e^{\Lambda t_i}\le e^{\Lambda t_q}/\Lambda$; and $\kappa_r\hat E(t_q)\le1$. Then apply Theorem~\ref{thm:allorders}.
\end{proof}

$\tau_{\rm nl}$ bounds the window over which the higher orders of the error stay below the first-order envelope $\hat E$; it says nothing about whether training on the window helps, which is the sign of $\langle g_v,g_w\rangle$ (Corollary~\ref{cor:sign}).
$C$, $\Lambda$, $\bar\eta$, $D_H$ and $r$ are properties of the trained network and must hold uniformly over starts and pairs of steps, so they are set by the strongest transient growth.
The factor $e^{\Lambda\dt}$ can equivalently be absorbed into $C$. If the argument of the logarithm is below one, $\tau_{\rm nl}$ is negative and the corollary gives no window.
For a direct model the same proof holds with $\bar\eta=\max_i|\eps_i|/\dt$ and $D_H$ replaced by $\max_iD_i/\dt$; whether $\max_iD_i$ is of order $\dt$ depends on the complex extension of the model's one-step map and is not implied by a small one-step error, and $C$, $\Lambda$ are not tied to $\dt$ either, since the learned Jacobian need not be close to the identity.

\paragraph{What these bounds cannot do.}
\begin{enumerate}
\item \emph{A converse.} A~priori norm bounds can certify that training on a window decreases the error at chosen horizons; they cannot prove that it fails to, because that needs control of cancellations in $\langle g_v,g_w\rangle$. The exact rate is computable a~posteriori, so failure can be checked numerically for a given network.
\item \emph{Tightness.} The triangle inequality is used at every step, so cancellations that reduce the error are invisible, and uniform constants are set by the worst start. No numerical comparison of \eqref{eq:gradbound} with the exact difference has been made.
\item \emph{Constants.} $r$, $D_i$, $D_i'$, $F_i$ are suprema over complex balls. None has been certified; the measured radius of Remark~\ref{rem:radius} is an upper bound on $r$. Certified values are possible by propagating through the layers that each SiLU unit stays holomorphic while the imaginary part of its pre-activation stays below $\pi-\delta$ in absolute value.
\end{enumerate}

\section{The local-complexity floor under multistep training}\label{sec:floor}

This section uses no expansion.

\paragraph{The roughness theory.}
The roughness theory concerns an affine method with exact lower slots, $G=Q_h+c_p\,\Nn$, where $Q_h$ is a known part, $\Nn$ the network slot and $c_p$ its slot weight; for an Euler model $Q_h=\mathrm{id}$, $c_p=\dt$, $\Nn=\Hop$, and for a direct model $Q_h=0$, $c_p=1$, $\Nn=\Gt$ (both have exact lower slots trivially).
The target of the network, in state units, is $U=S_h-Q_h$ with $S_h$ the true flow map over $h=\dt$ (written $\Phi_h$ in the roughness paper).
Its local complexity under a seminorm $\Rgh$ on sampled functions at the training inputs is $\Rgh(U)$, with norm constant $C_\Rgh$.
The regularity of the roughness paper's Assumption~1 (the inputs lie on a compact invariant set $\mathcal A$ of the flow on which the time derivatives $x^{(k)}$ are bounded, $M_k=\sup_{\mathcal A}|x^{(k)}|$, and Lipschitz) is used only for $\|S_h-\mathrm{id}\|\le M_1h$ in Corollary~\ref{cor:directrollout}.
The injected error at the training inputs is $E=G-S_h=c_p\Nn-U$; for noise-free samples at spacing $\dt$, $E$ at $\uu_T(t_i)$ is $-\eps_i$.
Let $\Psi_{\mathcal N}=\nabla_\theta\Nn\,\Pi^{1/2}$ at the training inputs, $K=\frac1{n_{\rm in}}\Psi_{\mathcal N}\Psi_{\mathcal N}^\top$ the training kernel ($n_{\rm in}$ the number of inputs) with eigenvalues $\lambda_i$ and eigenvectors $v_i$, $\sigma_\Rgh$ its peak feature roughness, $E_\gamma=\gamma(K+\gamma I)^{-1}E$ the residual at scale $\gamma$ and $\Upsilon(\gamma)=\sum_i\lambda_i\hat e_i^2/(\lambda_i+\gamma)^2$ its step size, $\hat e_i=\langle v_i,E\rangle$.
The roughness paper's theorem that local complexity bounds the injected error from below states
\begin{equation}\label{eq:floor}
\|E_\gamma\|\ \ge\ \frac1{C_\Rgh}\Big(\Rgh(U)-\Rgh(c_p\Nn)-\sigma_\Rgh\,\Upsilon(\gamma)^{1/2}\Big)\qquad\text{for every }\gamma>0,
\end{equation}
at every time during training, for any labels and any update rule.
For a single seminorm it is informative only while $\Rgh(c_p\Nn)+\sigma_\Rgh\Upsilon(\gamma)^{1/2}<\Rgh(U)$.

\paragraph{The best bound is the supremum over all roughness functions.}
By the roughness paper's proposition that the floor is a one-step bound, for $s_\gamma=\Upsilon(\gamma)^{1/2}/|c_p|$ every seminorm satisfies (right side of \eqref{eq:floor}) $\le b_\Rgh(s_\gamma)\le F(s_\gamma)=\|E_\gamma\|$, and by its supremum theorem the supremum over all seminorms equals $\|E_\gamma\|$, attained by the functional along $E_\gamma$.
By its Krylov theorem, the Krylov values $F_M(s_\gamma)$, built from $E,KE,\dots,K^{M-1}E$, increase to $\|E_\gamma\|$ and reach it after finitely many terms.
The supremum, certified by the Krylov values, is therefore the bound used below: for every $\gamma>0$ and $M$,
\begin{equation}\label{eq:supfloor}
\|E\|\ \ge\ \|E_\gamma\|\ \ge\ F_M(s_\gamma),
\end{equation}
and every bound \eqref{eq:floor} for a fixed seminorm is weaker. A fixed seminorm is needed only to express a bound through the data alone, before training, through $\Rgh(U)$.

\paragraph{Multistep training does not change the method.}
Training on a window loss changes how the parameters move, not the form of the step map: the known part, slot weight, network slot and target are those of one-step training.
For a window, take as training inputs its true states $\uu_T(t_i)$, all starts, $i<p$; this set, and with it $\Rgh$, $C_\Rgh$ and $\Rgh(U)$, changes with the window.

\begin{proposition}[The floor holds during multistep training]\label{prop:floorms}
At every time during training on a window loss, with the window's true states as inputs, the injected error satisfies \eqref{eq:floor} and \eqref{eq:supfloor}, with $K$, $\sigma_\Rgh$ and $\Rgh(c_p\Nn)$ those of the network at that time.
For a fixed set of inputs, $\Rgh(U)$ does not depend on the loss used for training.
\end{proposition}

\begin{proof}
The roughness paper's theorem is an inequality between quantities at one training time and holds for any update rule, in particular for the updates of multistep training.
\end{proof}

Here $K$ is the one-step kernel of the network slot at the inputs, not the kernel that drives multistep training ($\Psih\Psih^\top W$); under multistep training $E$ does not evolve by $-KE$, so $E_\gamma$ and the leak term act only as a spectral filter, without the roughness paper's reading as what gradient flow removes quickly.

\paragraph{From injected error to rollout error.}
For start $k$ and $i<p$ let $\ell_{k,i}=\sup\|\nabla\Gt\|$ over the segment joining $\uu_T(t_i)$ and $\uu_p(t_i)$, and $\ell=\max_{k,i}\ell_{k,i}$; $\ell$ is a~posteriori, and the segments can leave $\mathcal A$.

\begin{lemma}[Rollout error controls injected error]\label{lem:inverse}
For every start and $i<p$,
\[
\eps_i=\ee(t_{i+1})-\big(\Gt(\uu_T(t_i))-\Gt(\uu_p(t_i))\big),\qquad |\eps_i|\le|\ee(t_{i+1})|+\ell_{k,i}\,|\ee(t_i)| .
\]
With $\|\ee\|^2=\frac1{np}\sum_k\sum_{q=1}^p|\ee_k(t_q)|^2$ and $\|\eps\|^2=\frac1{np}\sum_k\sum_{i<p}|\eps_{k,i}|^2$,
\begin{equation}\label{eq:inverse}
\|\eps\|\le(1+\ell)\,\|\ee\| ,
\end{equation}
and for the time-averaged loss ($w_q=1/p$), $L_w=\tfrac12\|\ee\|^2\ge\tfrac12\|\eps\|^2/(1+\ell)^2$.
\end{lemma}

\begin{proof}
$\ee(t_{i+1})=\eps_i+\Gt(\uu_T(t_i))-\Gt(\uu_p(t_i))$, and the mean value inequality bounds the difference by $\ell_{k,i}|\ee(t_i)|$.
By Cauchy--Schwarz, $(a+\ell b)^2\le(1+\ell)(a^2+\ell b^2)$; summing over $i<p$ with $\ee(t_0)=0$, $\sum_{i<p}|\eps_i|^2\le(1+\ell)\big(\sum_{q=1}^p|\ee(t_q)|^2+\ell\sum_{i=1}^{p-1}|\ee(t_i)|^2\big)\le(1+\ell)^2\sum_{q=1}^p|\ee(t_q)|^2$.
\end{proof}

\begin{theorem}[Local-complexity floor for the rollout error]\label{thm:rolloutfloor}
At every time during training, for any loss and any update rule, with all quantities of \eqref{eq:floor} built on the window's true states and the same normalization, the window-averaged root-mean-square rollout error satisfies
\begin{equation}\label{eq:rolloutfloor}
\|\ee\|\ \ge\ \frac{\|E\|}{1+\ell}\ \ge\ \frac{\|E_\gamma\|}{1+\ell}\ \ge\ \frac{F_M(s_\gamma)}{1+\ell}\qquad\text{for every }\gamma>0\text{ and }M,
\end{equation}
and, for any fixed seminorm, the weaker $\|\ee\|\ge\big(\Rgh(U)-\Rgh(c_p\Nn)-\sigma_\Rgh\Upsilon(\gamma)^{1/2}\big)/\big((1+\ell)C_\Rgh\big)$.
\end{theorem}

\begin{proof}
$\|E\|=\|\eps\|$, Lemma~\ref{lem:inverse}, \eqref{eq:supfloor} and \eqref{eq:floor}.
\end{proof}

For Euler models $\ell\le1+\dt\sup\|\nabla\Hop\|$ over the segments; the factor $1+\ell$ is close to $2$ when $\dt\|\nabla\Hop\|\ll1$ there.

\begin{corollary}[Direct models: the unbuilt identity]\label{cor:directrollout}
For a direct model at any training time let $u=\mathrm{id}-\Gt$ be the identity residual at the inputs and $F_u(s)=\min_{|g|\le s}\|u-\Psi_{\mathcal N}g\|$, the supremum over all seminorms of the corresponding step bounds, certified by its Krylov values (the roughness paper's corollary for any training time).
Every step of length at most $s$ leaves an injected error at least $F_u(s)-\|S_h-\mathrm{id}\|\ge F_u(s)-M_1h$ (with the step-averaged features this bounds the step actually taken exactly; with the current features it is first order in $s$), so by Lemma~\ref{lem:inverse} the rollout error over the window after that step satisfies
\[
\|\ee_{\rm new}\|\ \ge\ \frac{F_u(s)-M_1h}{1+\ell_{\rm new}} ,
\]
with $\ell_{\rm new}$ evaluated after the step.
\end{corollary}

$F_u(s)$, the part of the identity map a step of length $s$ cannot build, does not depend on $h$ except through the network, while the subtracted term is of order $h$; the bound is positive only while the direct network has not built the identity to within $M_1h$ by steps of length $s$, and it holds under any training loss.
For an Euler model, by the roughness paper's comparison of the two steps, the smallest error after a step of any length is at most $\|S_h-\mathrm{id}\|\le M_1h$.

What can move the floor is the network's current output and its kernel at the inputs (through $E_\gamma$), the factor $1+\ell$ and the choice of inputs; the target is fixed by the inputs and the method.

\paragraph{The kernel-level results.}
The roughness paper's rates at one training time (continuous time) and its bound for every step of length at most $s$ use only a linear map from parameter steps to sampled functions; with the rollout features $\Psih$, the rollout error $\ee$, the inner product $\langle a,b\rangle_w=\frac1n\sum_{k,q}w_qa_{k,q}\cdot b_{k,q}$ (restricted to horizons with $w_q>0$) and seminorms on functions of start and horizon, they hold exactly for multistep training; in the every-step bound the maxima over the parameter ball are of the rollout features $\Psih(\theta')$, which include the change of the predicted trajectory.
With the rollout features averaged along the step actually taken, $\min_{|g'|\le s}\|\ee-\bar\Psi g'\|_w$ is an exact lower bound on the error after that step (first statement of the one-step-training theorem of the rollout-kernel note).
As statements about the best of all steps of length $s$, the smallest error after one step, its slow-part formula, attainment by a rank-one functional and the Krylov series hold exactly only through the every-step bound, and otherwise only to first order in $s$, with no bounded remainder; rollout features can change quickly with the parameters at long windows.
None of these involves the target $U$ or its local complexity.

\section{Errors of the learned operators along the trajectory}\label{sec:operators}

The coefficients of the error expansion are the state derivatives $\nabla^k\Gt$ of the learned map along the true trajectory.
This section splits each of them into the corresponding derivative of the true flow map and an error of the learned operator, and applies the roughness theory to each operator error.
The parameter features $\Phi$ and the mixed derivatives $\nabla^k_\uu\nabla_\theta\Gt$ in the gradient and rate expansions have no counterpart in the true system and are not split.
In this section $k$ denotes the order of a derivative, not a start.

\paragraph{Model error and operator errors.}
Let $E(\uu)=\Gt(\uu)-S_h(\uu)$ be the model error as a function on state space; at the data, $E(\uu_T(t_i))=-\eps_i$.
For $k\ge1$ let $\Delta_k(\uu)=\nabla^kE(\uu)=\nabla^k\Gt(\uu)-\nabla^kS_h(\uu)$, the error of the $k$-th learned operator, and $\Delta_0=E$.
Along the true trajectory write $\Delta_{k,i}=\Delta_k(\uu_T(t_i))$, $J^S_i=\nabla S_h(\uu_T(t_i))$, so $J_i=J^S_i+\Delta_{1,i}$, and let $\Pp_S(t_q,t_i)=J^S_{q-1}\cdots J^S_{i+1}$.
By the chain rule $\Pp_S(t_q,t_i)=\nabla S_{t_q-t_{i+1}}(\uu_T(t_{i+1}))$: the true propagators are the tangent flow of the true system along the data.

\begin{lemma}[Rollout error: true separation plus model error at the predicted state]\label{lem:split}
Exactly, for every start and $i<p$,
\[
\ee(t_{i+1})=\big(S_h(\uu_T(t_i))-S_h(\uu_p(t_i))\big)-E(\uu_p(t_i)).
\]
If $S_h$ and $E$ are analytic and $|\ee(t_i)|$ is within the radii of their Taylor series at $\uu_T(t_i)$ in the direction of $\ee(t_i)$, then
\[
\ee(t_{i+1})=\big(J^S_i+\Delta_{1,i}\big)\ee(t_i)+\eps_i+\sum_{k\ge2}\frac{(-1)^{k+1}}{k!}\big(\nabla^kS_h(\uu_T(t_i))+\Delta_{k,i}\big)\big[\ee(t_i)^{\otimes k}\big].
\]
\end{lemma}

\begin{proof}
$\uu_T(t_{i+1})=S_h(\uu_T(t_i))$ and $\Gt=S_h+E$, so $\ee(t_{i+1})=S_h(\uu_T(t_i))-S_h(\uu_p(t_i))-E(\uu_p(t_i))$. Expand both terms about $\uu_T(t_i)$ with $\uu_p(t_i)=\uu_T(t_i)-\ee(t_i)$ and use $E(\uu_T(t_i))=-\eps_i$.
\end{proof}

The rollout error is the separation of the true and predicted states by the true flow, plus the model error evaluated off the data, at the predicted state.
Its coefficients are those of Lemma~\ref{lem:recursion}, split as $\nabla^k\Gt=\nabla^kS_h+\Delta_k$; at each order in $\eps$ the error expansion of Section~\ref{sec:error} therefore becomes a polynomial in $(\Delta_1,\Delta_2,\dots)$ (with $\Delta_1$ entering through every propagator) whose coefficients are derivatives of the true flow map; the gradient and rate expansions contain in addition $\Phi$ and its mixed derivatives, which are not split.
Under Assumption~\ref{ass:analytic}, $g_i$ splits in the same way, so $D_i\le D^S_i+D^E_i$ and $D'_i\le D'^S_i+D'^E_i$, with the superscripts denoting the corresponding suprema for $S_h(\uu_i+z)-S_h(\uu_i)-z$ and for $E(\uu_i+z)-E(\uu_i)$, whenever both extend; $F_i$ has no true counterpart.

\begin{proposition}[Learned and true propagators]\label{prop:propsplit}
Exactly,
\[
\Pp(t_q,t_i)-\Pp_S(t_q,t_i)=\sum_{i<j<q}\Pp_S(t_q,t_j)\,\Delta_{1,j}\,\Pp(t_j,t_i).
\]
If $\|\Pp_S(t_j,t_i)\|\le C_Se^{\Lambda_S(t_j-t_i)}$ for all $i<j$ ($C_S\ge1$), then
\begin{align*}
\|\Pp(t_q,t_i)\|&\le C_Se^{\Lambda_S(t_q-t_i)}\prod_{i<j<q}\big(1+C_S\|\Delta_{1,j}\|\big),\\
\|\Pp(t_q,t_i)-\Pp_S(t_q,t_i)\|&\le C_Se^{\Lambda_S(t_q-t_i)}\Big(\prod_{i<j<q}\big(1+C_S\|\Delta_{1,j}\|\big)-1\Big)\le C_Se^{\Lambda_S(t_q-t_i)}\Big(e^{C_S\sum_{i<j<q}\|\Delta_{1,j}\|}-1\Big),
\end{align*}
and the first term of the error satisfies $\big|\ee^{(1)}(t_q)-\sum_{i<q}\Pp_S(t_q,t_i)\eps_i\big|\le\sum_{i<q}\|\Pp(t_q,t_i)-\Pp_S(t_q,t_i)\|\,|\eps_i|$.
\end{proposition}

\begin{proof}
The identity is the telescoping of $J_{q-1}\cdots J_{i+1}-J^S_{q-1}\cdots J^S_{i+1}$ with $J_j-J^S_j=\Delta_{1,j}$.
With $z_q=\|\Pp(t_q,t_i)\|e^{-\Lambda_S(t_q-t_i)}$ it gives $z_q\le C_S+\sum_{i<j<q}C_S\|\Delta_{1,j}\|z_j$, and $z_{i+1}=e^{-\Lambda_S\dt}\le C_S$ by the hypothesis at $j=i+1$ ($1=\|\I\|\le C_Se^{\Lambda_S\dt}$), so the discrete Grönwall inequality gives the first bound.
Inserting it into the identity and using $\sum_ja_j\prod_{l<j}(1+a_l)=\prod_j(1+a_j)-1$ gives the second, and $1+a\le e^a$ the last.
\end{proof}

$C_S$ and $\Lambda_S$ belong to the true system along the data; the network enters only through the Jacobian errors. The proposition concerns the propagators along the true trajectory; those along the predicted trajectory, $\Ph$, which enter Theorem~\ref{thm:gradbound} and $B(\rho)$, are not covered.
For an Euler model $\Delta_{1,j}=\dt\,(\nabla\Hop-\nabla f_h)(\uu_T(t_j))$ with $f_h=(S_h-\mathrm{id})/\dt$, so $\sum_{i<j<q}\|\Delta_{1,j}\|\le(t_q-t_i)\sup_j\|\nabla\Hop-\nabla f_h\|$: the difference between learned and true propagators stays below a fixed multiple of the bound $C_Se^{\Lambda_S(t_q-t_i)}$ over windows of length up to a constant times $1/(C_S\sup\|\nabla\Hop-\nabla f_h\|)$.
This scale is set by the Jacobian error, and is separate from the logarithmic dependence on the injection error of Corollary~\ref{cor:closed}.

\paragraph{The roughness theory for every operator.}
For $k\ge0$ the $k$-th operator of the method is $\nabla^kG=\nabla^kQ_h+c_p\nabla^k\Nn$, with target $\nabla^kU=\nabla^kS_h-\nabla^kQ_h$, fixed by the system and the method, and error $\Delta_k=c_p\nabla^k\Nn-\nabla^kU$ at the training inputs, a sampled tensor-valued function.
Tensor-valued sampled functions carry the inner product $\langle a,b\rangle=\frac1{n_{\rm in}}\sum_i\langle a_i,b_i\rangle_F$ (Frobenius at each input) and its norm; seminorms, norm constants and peak feature roughness are taken with respect to it.
The gradient features of the $k$-th operator are $\Xi^{(k)}=\partial_\theta\nabla^k\Nn\,\Pi^{1/2}=\nabla_\uu^k\nabla_\theta\Nn\,\Pi^{1/2}$ at the inputs (mixed partial derivatives commute), with kernel $K_k=\frac1{n_{\rm in}}\Xi^{(k)}\Xi^{(k)\top}$, residual $\Delta_{k,\gamma}=\gamma(K_k+\gamma I)^{-1}\Delta_k$, step size $\Upsilon_k(\gamma)$, $s^{(k)}_\gamma=\Upsilon_k(\gamma)^{1/2}/|c_p|$, peak feature roughness $\sigma^{(k)}_\Rgh$, single-step value $F^{(k)}(s)=\min_{|\mathbf g|\le s}\|\Delta_k+c_p\Xi^{(k)}\mathbf g\|$ and Krylov values $F^{(k)}_M$ built from $\Delta_k,K_k\Delta_k,\dots$; $k=0$ is Section~\ref{sec:floor}.

\begin{theorem}[Each operator error obeys the roughness theory]\label{thm:operators}
At every training time, for any loss and any update rule:
\begin{enumerate}
\item $\dot\Delta_k=c_p\Xi^{(k)}\hat g$, so $|\tfrac{d}{dt}\langle\varphi,\Delta_k\rangle|\le|c_p|\,\|\varphi\|_{K_k}|\hat g|$ for every sampled $\varphi$.
\item For every seminorm, $\|\Delta_{k,\gamma}\|\ge\big(\Rgh(\nabla^kU)-\Rgh(c_p\nabla^k\Nn)-\sigma^{(k)}_\Rgh\Upsilon_k(\gamma)^{1/2}\big)/C_\Rgh$; the supremum over seminorms is $\|\Delta_{k,\gamma}\|=F^{(k)}(s^{(k)}_\gamma)$, certified by its Krylov values: $\|\Delta_k\|\ge\|\Delta_{k,\gamma}\|\ge F^{(k)}_M(s^{(k)}_\gamma)$ for every $\gamma>0$ and $M$.
\item With the features averaged along a step, the analogue $\bar F^{(k)}(s)$ of $F^{(k)}(s)$ bounds that step's $k$-th operator error from below exactly; for every step of length at most $s$, the roughness paper's every-step bound with the features $\Xi^{(k)}(\theta')$ over the parameter ball gives an exact lower bound; $F^{(k)}(s)$ as the best value reachable by any step of length $s$ holds only to first order in $s$.
\end{enumerate}
Under gradient flow on a loss $X$, $\dot\Delta_k=-c_p\Xi^{(k)}g_X$. For the one-step loss $\tfrac12\|\Nn-U/c_p\|^2$ this is $\dot\Delta_k=-K_{k0}E$ with the cross-kernel $K_{k0}=\frac1{n_{\rm in}}\Xi^{(k)}\Xi^{(0)\top}$; for the window loss it is $\dot\Delta_k=c_p\Xi^{(k)}\Psih^\top W\ee$.
\end{theorem}

\begin{proof}
$\Delta_k=c_p\nabla^k\Nn-\nabla^kU$ has the form ``slot weight times a parametrized output minus a fixed target'', and $\partial_\theta(c_p\nabla^k\Nn)\Pi^{1/2}=c_p\Xi^{(k)}$; the roughness paper's rate proposition, floor theorem, proposition that the floor is a one-step bound, supremum theorem, Krylov theorem and every-step bound are finite-dimensional Hilbert-space statements that use nothing else.
The last statements are $\hat g=-g_X$ with $g_X=\frac1{n_{\rm in}c_p}\Xi^{(0)\top}E$ for the one-step loss, and $g_w=-\Psih^\top W\ee$.
\end{proof}

For $k\ge1$, even under one-step training $\Delta_k$ does not evolve by $-K_k\Delta_k$, so $\Delta_{k,\gamma}$ and the leak term act only as a spectral filter, as for $k=0$ under multistep training.
One-step training is chosen from the values alone and reaches the $k$-th operator only through the cross-kernel $K_{k0}$; the window loss reaches it through $\Xi^{(k)}\Psih^\top$, whose rollout features contain the mixed derivatives of Section~\ref{sec:grad}.
For $k=1$, $F^{(1)}_M$ certifies a part of the Jacobian error that a step of length $s$ cannot remove, and Proposition~\ref{prop:propsplit} turns Jacobian errors into propagator errors.

\begin{proposition}[Local complexity of the operator targets]\label{prop:opcomplexity}
Assume the roughness paper's Assumption~1, that the inputs lie in $\mathcal A$, that $(t,x)\mapsto S_t(x)$ is smooth, and let $\bar M^{(k)}_{j}=\sup_{x\in\mathcal A,\,0\le t\le h}\|\nabla_x^k[x^{(j)}(S_t(x))]\|_F$.
For the Taylor method of order $\bar p\ge1$ with exact lower slots and every seminorm,
\[
\Big|\Rgh(\nabla^kU)-\frac{h^{\bar p}}{\bar p!}\Rgh\big(\nabla^kx^{(\bar p)}\big)\Big|\le\frac{C_\Rgh\bar M^{(k)}_{\bar p+1}h^{\bar p+1}}{(\bar p+1)!},
\]
and for the direct step $|\Rgh(\nabla^kU)-\Rgh(\nabla^k\mathrm{id})|\le C_\Rgh\bar M^{(k)}_1h$, with $\nabla\mathrm{id}=\I$ and $\nabla^k\mathrm{id}=0$ for $k\ge2$.
\end{proposition}

\begin{proof}
$\partial_t^jS_t(x)=x^{(j)}(S_t(x))$, so $\partial_t^j\nabla_x^kS_t(x)|_{t=0}=\nabla^kx^{(j)}(x)$ and $\partial_t^{\bar p+1}\nabla_x^kS_t(x)=\nabla_x^k[x^{(\bar p+1)}(S_t(x))]$; the chain rule through $S_t$ is contained in $\bar M^{(k)}_j$, which is finite by compactness of $\mathcal A\times[0,h]$.
Taylor's formula in $t$ with integral remainder gives $\nabla^kU=\nabla^kS_h-\sum_{j<\bar p}\frac{h^j}{j!}\nabla^kx^{(j)}=\frac{h^{\bar p}}{\bar p!}\nabla^kx^{(\bar p)}+R$ with $\sup_{\mathcal A}\|R\|_F\le\bar M^{(k)}_{\bar p+1}h^{\bar p+1}/(\bar p+1)!$, and $|\Rgh(a+R)-\Rgh(a)|\le\Rgh(R)\le C_\Rgh\|R\|\le C_\Rgh\sup_{\mathcal A}\|R\|_F$ because the inputs lie in $\mathcal A$. The direct step is $\bar p=0$ with $Q_h=0$.
\end{proof}

For an Euler model the Jacobian target is $h\nabla f+O(h^2)$, of the same order in $h$ as its value target; for a direct model it is $\I+O(h)$, a constant plus a term of order $h$ (seminorms that vanish on constants give it local complexity of order $h$), and the targets of the operators of order $k\ge2$ are of order $h$ for both.

\paragraph{The size of the injection errors does not decide whether a window helps.}
By Proposition~\ref{prop:rateorders} the leading term of the rate is a quadratic form in the injections, $R^{(2)}=\eps^\top A\eps$ with $A=\LL^\top W_v\LL\,\Phi\Phi^\top\LL^\top W\LL$.

\begin{proposition}[Norms do not determine the sign]\label{prop:normsign}
Let Assumption~\ref{ass:analytic} hold, hold all coefficients (the true trajectories, the derivatives of $\Gt$ along them and $\Phi$) fixed, and replace the injections by hypothetical injections $s\hat\eps$, as in the scaling of Section~\ref{sec:analytic}.
If the symmetric part of $A$ has eigenvalues of both signs (which needs $v\ne w$), there are $\hat\eps_\pm$ with $|\hat\eps_+|=|\hat\eps_-|$ such that for all sufficiently small $s>0$ training on $L_w$ decreases $V_v$ with injections $s\hat\eps_+$ and increases it with injections $s\hat\eps_-$.
Hence, as a function of the injections with the coefficients held fixed, no criterion that depends on the injections only through their total norm determines the sign of the rate for small injections.
\end{proposition}

\begin{proof}
Take $\hat\eps_\pm$ unit eigenvectors of $\tfrac12(A+A^\top)$ with eigenvalues of opposite signs, so $R^{(2)}(\hat\eps_+)>0>R^{(2)}(\hat\eps_-)$.
Fix $\rho=2$ and compute the sequence $a$ of Theorem~\ref{thm:allorders} for the injections $s\hat\eps_\pm$; as $s\to0$, $\rho s|\LL\hat\eps_\pm|+a(\rho s)\to0$, so $\rho$ is admissible for small $s$. By Theorem~\ref{thm:tradeoffbound} the rate is holomorphic in $s$ with $-\dot V_v=s^2R^{(2)}(\hat\eps)+O(s^3)$, and its sign is that of $R^{(2)}(\hat\eps)$ for small $s$.
\end{proof}

No network realizing the injections $s\hat\eps_\pm$ with these coefficients is exhibited; the statement concerns what can be inferred from the norm of the injections alone.

What is needed instead is the structure of the injections and operator errors relative to the propagators and the kernel: their joint moments.

\paragraph{The window for one training step, in expectation.}
Fix target horizons $v$ and a preconditioned gradient step of size $\alpha$ on window $w$, $\Delta\theta=-\alpha\Pi\nabla_\theta L_w$ (rescaled length $\alpha|g_w|$); $V_v$ and $L_w$ use the same starts.
With $H_v=\Pi^{1/2}\nabla^2_\theta V_v\Pi^{1/2}$, Taylor's formula with integral remainder gives exactly
\begin{equation}\label{eq:onestepdecrease}
V_v(\theta+\Delta\theta)-V_v(\theta)=-\alpha\langle g_v,g_w\rangle+\alpha^2\int_0^1(1-t)\,g_w^\top H_v\big(\theta+t\Delta\theta\big)\,g_w\,dt .
\end{equation}
Let the starts be drawn from a distribution (the data), write $\mathbb E$ for the expectation over draws, and let $\mathcal G_v(w,\alpha)$ be the expected decrease, minus the expectation of \eqref{eq:onestepdecrease}.
The window for this step is a maximizer of $\mathcal G_v(w,\alpha_w)$ over $w$, with $\alpha_w$ the step size used for window $w$ (for example from the local quadratic model along $g_w$).
The expected rate $\mathbb E\langle g_v,g_w\rangle$ is the first-order part of $\mathcal G_v$.

\begin{proposition}[Expected rate to all orders]\label{prop:expectedrate}
Each $R^{(M)}$ of \eqref{eq:RM} is a sum of products of the injections $\eps$, the operator errors $\Delta_1,\dots,\Delta_{M-1}$, the derivatives of the true flow map and the parameter features $\Phi$ and their mixed derivatives, all along the true trajectories; since all of these vary with the start, $\mathbb E\,R^{(M)}$ is a joint moment of these quantities, across the steps of a trajectory and across the pairs of starts coupled by the kernel.
Suppose that for a fixed $\rho>1$, $\rho$ is admissible for almost every draw of starts and $\mathbb E[B(\rho)]<\infty$, with $B(\rho)$ of Theorem~\ref{thm:tradeoffbound}.
Then $\mathbb E\langle g_v,g_w\rangle=\sum_{M\ge2}\mathbb E\,R^{(M)}$ and $\big|\mathbb E\langle g_v,g_w\rangle-\sum_{M=2}^N\mathbb E\,R^{(M)}\big|\le\mathbb E[B(\rho)]\rho^{-N}/(\rho-1)$.
\end{proposition}

\begin{proof}
The first statement follows from \eqref{eq:orders}, \eqref{eq:Phij}--\eqref{eq:Psij} with $\nabla^k\Gt=\nabla^kS_h+\Delta_k$, by multilinearity.
For each draw with $\rho$ admissible, Theorem~\ref{thm:tradeoffbound} gives $\langle g_v,g_w\rangle=\sum_MR^{(M)}$ and $|R^{(M)}|\le B(\rho)\rho^{-M}$; hence $\sum_M\mathbb E|R^{(M)}|\le\mathbb E[B(\rho)]/(\rho-1)<\infty$, Fubini's theorem exchanges expectation and sum, and the tail bound gives the remainder.
\end{proof}

The second term of \eqref{eq:onestepdecrease} involves $H_v$ along the step, where all coefficients change; Section~\ref{sec:stepseries} treats the step exactly as a series in the step size.
The moments are an input: they come from the data, or from a model of the injections and operator errors; Proposition~\ref{prop:normsign} shows that a model through the total norm alone is not enough.
We expect the admissibility condition to restrict the window: $B(\rho)$ grows with the propagators of each start, so at long windows the few fastest-stretching starts can fail to be admissible, and then only a truncation with the remainder on the admissible draws is available.
The choice is for one training step at the current parameters.

\section{One training step: expansion in the step to any order}\label{sec:stepseries}

The previous sections expand in the one-step errors at fixed parameters.
This section expands in the training step itself, for a fixed window of $p$ physical steps, to any order, with a bounded remainder.
It describes how the window error responds to a single step of a given length, and how that response changes with $p$.

\paragraph{Setting.}
Fix the parameters $\theta$, a window of $p$ steps with weights $w_q\ge0$, $q\le p$, the window error $L_w=\tfrac12\ee^\top W\ee$ over horizons $1,\dots,p$ (with the average over the $n$ starts in $W$), its rescaled gradient $g_w=-\Psih^\top W\ee$, and a unit rescaled direction $\mathbf d$, $|\mathbf d|=1$ (for example $\mathbf d=-g_w/|g_w|$).
For a scalar $z$ let $\uu_q(z)$ be the rollout with parameters $\theta+z\,\Pi^{1/2}\mathbf d$ from the same start, $\uu_0(z)=\uu_T(t_0)$, so that the step has rescaled length $|z|$; let $\ee_q(z)=\uu_T(t_q)-\uu_q(z)$ and $L_w(z)=\tfrac12\ee(z)^\top W\ee(z)$.
Write $\mathbf Y^{(0)}_q=\uu_p(t_q)$ and $\mathbf Y^{(m)}_q=\partial_z^m\uu_q(z)|_{z=0}$, $m\ge1$, for the sensitivities of order $m$ along $\mathbf d$; $\mathbf Y^{(1)}_q=\Psih_q\mathbf d$.
Let $D^k\Gt$ be the $k$-th derivative of $\Gt$ jointly in state and parameters, and $\mathbf w^{(1)}_q=(\mathbf Y^{(1)}_q,\Pi^{1/2}\mathbf d)$, $\mathbf w^{(j)}_q=(\mathbf Y^{(j)}_q,0)$ for $j\ge2$, the derivatives of the curve $z\mapsto(\uu_q(z),\theta+z\Pi^{1/2}\mathbf d)$.
All constants below depend on $\mathbf d$ (and, for $\mathbf d=-g_w/|g_w|$, on the window).

\begin{lemma}[Sensitivities of every order]\label{lem:sens}
For every $m\ge1$, $\mathbf Y^{(m)}_0=0$ and
\[
\mathbf Y^{(m)}_{q+1}=\hat J_q\,\mathbf Y^{(m)}_q+\mathbf S^{(m)}_q,\qquad
\mathbf S^{(m)}_q=\sum_{\pi\in\mathcal P_m,\ |\pi|\ge2}D^{|\pi|}\Gt\big(\uu_p(t_q);\theta\big)\big[\mathbf w^{(|B|)}_q\big]_{B\in\pi}+[m=1]\,\Phih_q\mathbf d,
\]
where $\mathcal P_m$ is the set of partitions of $\{1,\dots,m\}$ into blocks $B$; hence $\mathbf Y^{(m)}_q=\sum_{i<q}\Ph(t_q,t_i)\mathbf S^{(m)}_i$.
The source $\mathbf S^{(m)}_q$ involves only $\mathbf Y^{(j)}_q$ with $j<m$.
For $m=2$, with all derivatives at $(\uu_p(t_q);\theta)$, $\mathbf S^{(2)}_q=\nabla^2_\uu\Gt[\mathbf Y^{(1)}_q,\mathbf Y^{(1)}_q]+2\nabla_\uu\nabla_\theta\Gt[\mathbf Y^{(1)}_q,\Pi^{1/2}\mathbf d]+\nabla^2_\theta\Gt[\Pi^{1/2}\mathbf d,\Pi^{1/2}\mathbf d]$.
\end{lemma}

\begin{proof}
Differentiate $\uu_{q+1}(z)=\Gt(\uu_q(z);\theta+z\Pi^{1/2}\mathbf d)$ $m$ times at $z=0$ with the multivariate Faà di Bruno formula, $\partial_z^mF(\mathbf w(z))=\sum_{\pi\in\mathcal P_m}D^{|\pi|}F[\mathbf w^{(|B|)}]_{B\in\pi}$. The partition with one block gives $D\Gt[\mathbf w^{(m)}_q]=\hat J_q\mathbf Y^{(m)}_q+[m=1]\Phih_q\mathbf d$, since the parameter part of the curve is linear in $z$. A partition with at least two blocks has blocks of size below $m$. Unroll from $\mathbf Y^{(m)}_0=0$.
\end{proof}

For $m=1$ this is the features lemma of the rollout-kernel note; every higher order is carried by the same propagators $\Ph$, with sources built from the derivatives of the learned map in state and parameters along the rollout.

\begin{assumption}[Analytic extension along the step]\label{ass:stepanalytic}
There are $r_u,r_\theta,L_u,L_\theta>0$ such that for every start and every $q<p$ the map $(\bm\xi,z)\mapsto\Gt(\bm\xi;\theta+z\Pi^{1/2}\mathbf d)$ extends holomorphically to $\mathcal D_q=\{|\bm\xi-\uu_p(t_q)|<r_u,\ |z|<r_\theta\}\subset\mathbb C^d\times\mathbb C$, with $\|\nabla_{\bm\xi}\Gt\|\le L_u$ and $|\partial_z\Gt|\le L_\theta$ on $\mathcal D_q$ (Hermitian norms and complex operator norms).
\end{assumption}

Let $\phi_q=\sum_{j<q}L_u^j$ (so $\phi_0=0$, $\phi_{q+1}=1+L_u\phi_q$) and
\[
R_p=\min\Big(r_\theta,\ \frac{r_u}{L_\theta\,\phi_p}\Big).
\]

\begin{lemma}[The rollout is holomorphic in the step]\label{lem:stepholo}
Under Assumption~\ref{ass:stepanalytic}, for $|z|<R_p$ the rollout $\uu_q(z)$, $q\le p$, is defined by the recursion with complex $z$, is holomorphic in $z$, and satisfies $|\uu_q(z)-\uu_q(0)|\le L_\theta\phi_q|z|$.
\end{lemma}

\begin{proof}
Induction on $q$, with $\bm\delta_q=\uu_q(z)-\uu_q(0)$ and $\bm\delta_0=0$. If $|\bm\delta_q|\le L_\theta\phi_q|z|<r_u$ and $|z|<r_\theta$, both $(\uu_q(0)+\bm\delta_q,z)$ and $(\uu_q(0),z)$ lie in $\mathcal D_q$, and
$\bm\delta_{q+1}=\big[\Gt(\uu_q(0)+\bm\delta_q;\theta+z\Pi^{1/2}\mathbf d)-\Gt(\uu_q(0);\theta+z\Pi^{1/2}\mathbf d)\big]+\big[\Gt(\uu_q(0);\theta+z\Pi^{1/2}\mathbf d)-\Gt(\uu_q(0);\theta)\big]$.
The mean value inequality (the integral of the derivative along each segment, which lies in the convex set $\mathcal D_q$) gives $|\bm\delta_{q+1}|\le L_u|\bm\delta_q|+L_\theta|z|\le L_\theta\phi_{q+1}|z|$, which is below $r_u$ because $\phi_{q+1}\le\phi_p$. Each step composes holomorphic maps.
\end{proof}

\begin{theorem}[The step series and its remainder]\label{thm:stepseries}
Under Assumption~\ref{ass:stepanalytic}, for $|z|<R_p$:
\begin{enumerate}
\item $\uu_q(z)=\sum_{m\ge0}\frac{z^m}{m!}\mathbf Y^{(m)}_q$, with $\mathbf Y^{(m)}$ from Lemma~\ref{lem:sens}, and for every $R\in(|z|,R_p)$
\[
\frac{|\mathbf Y^{(m)}_q|}{m!}\le L_\theta\phi_qR^{1-m}\quad(m\ge1),\qquad
\Big|\uu_q(z)-\sum_{m=0}^M\frac{z^m}{m!}\mathbf Y^{(m)}_q\Big|\le L_\theta\phi_qR\,\frac{(|z|/R)^{M+1}}{1-|z|/R}.
\]
\item $L_w(z)=\sum_{m\ge0}\ell_mz^m$ with
\[
\ell_0=L_w,\qquad \ell_m=-\Big\langle\ee,\frac{\mathbf Y^{(m)}}{m!}\Big\rangle_W+\frac12\sum_{\substack{a+b=m\\ a,b\ge1}}\Big\langle\frac{\mathbf Y^{(a)}}{a!},\frac{\mathbf Y^{(b)}}{b!}\Big\rangle_W,
\]
where $\langle a,b\rangle_W=a^\top Wb$; in particular $\ell_1=\langle g_w,\mathbf d\rangle$ and $\ell_2=\tfrac12\|\Psih\mathbf d\|_W^2-\tfrac12\langle\ee,\mathbf Y^{(2)}\rangle_W$, half the curvature of $L_w$ along $\mathbf d$, split into its Gauss--Newton part and the error-weighted part (Section~\ref{sec:curvature}).
\item For every $R\in(|z|,R_p)$ and $M\ge0$,
\[
\Big|L_w(z)-\sum_{m=0}^M\ell_mz^m\Big|\le B_L(R)\frac{(|z|/R)^{M+1}}{1-|z|/R},\qquad B_L(R)=\|\ee\|_W\,L_\theta R\,\|\phi\|_W+\tfrac12L_\theta^2R^2\|\phi\|_W^2,
\]
with $\|\phi\|_W^2=\sum_qw_q\phi_q^2$ (the average over starts in $W$ included; $\phi_q$ is the same for all starts because the constants are uniform).
\end{enumerate}
\end{theorem}

\begin{proof}
By Lemma~\ref{lem:stepholo}, $\uu_q$ is holomorphic on $|z|<R_p$; its Taylor coefficients at $0$ are $\mathbf Y^{(m)}_q/m!$, and on $|z|=R$, $|\uu_q(z)-\uu_q(0)|\le L_\theta\phi_qR$, so Cauchy's estimate gives the coefficient bound for $m\ge1$ and summing the tail gives the remainder.
$L_w(z)$, defined with the bilinear form (no complex conjugation), is holomorphic on the same disc; inserting $\ee(z)=\ee-\bm\delta(z)$ with $\bm\delta(z)=\sum_{m\ge1}\frac{z^m}{m!}\mathbf Y^{(m)}$ gives $\ell_m$; $\ell_1=-\ee^\top W\Psih\mathbf d=\langle g_w,\mathbf d\rangle$.
On $|z|=R$, $|L_w(z)-\ell_0|=|{-\ee^\top W\bm\delta(z)}+\tfrac12\bm\delta(z)^\top W\bm\delta(z)|\le\|\ee\|_W\|\bm\delta(z)\|_W+\tfrac12\|\bm\delta(z)\|_W^2\le B_L(R)$ by Cauchy--Schwarz for the Hermitian seminorm of the diagonal $W\ge0$, so $|\ell_m|\le B_L(R)R^{-m}$ for $m\ge1$, and summing over $m>M$ gives the remainder.
\end{proof}

\paragraph{The single step and the roughness theory.}
For each direction, $L_w(z\mathbf d)$ differs from the model $\tfrac12\|\ee-z\Psih\mathbf d\|_W^2$, in which the error is linearized in the step with the current features, by $-\tfrac12\langle\ee,\mathbf Y^{(2)}\rangle_Wz^2+\sum_{m\ge3}\ell_mz^m$.
The roughness paper's single-step value for the window is the minimum of that model over all steps, $\tfrac12F_w(s)^2=\min_{|\mathbf d|=1,\,0\le z\le s}\tfrac12\|\ee-z\Psih\mathbf d\|_W^2$, certified from below by its Krylov values.
Theorem~\ref{thm:stepseries} bounds the difference along any given direction; a statement comparable to $F_w(s)$ for the best step over all directions needs the constants of Assumption~\ref{ass:stepanalytic} uniform over $|\mathbf d|=1$, and then gives $\min_{\mathbf d}L_w(s\mathbf d)\ge\tfrac12F_w(s)^2-\sup_{\mathbf d}\big|\tfrac12\langle\ee,\mathbf Y^{(2)}\rangle_Ws^2-\sum_{m\ge3}\ell_ms^m\big|$.

\begin{corollary}[Certified decrease of the window error by one step]\label{cor:stepcap}
Under Assumption~\ref{ass:stepanalytic}, a step of rescaled length $s<R_p$ along $\mathbf d$ decreases the window error, $L_w(s)<L_w$, if for some $M\ge1$ and $R\in(s,R_p)$
\[
\sum_{m=1}^M\ell_ms^m+B_L(R)\frac{(s/R)^{M+1}}{1-s/R}<0 .
\]
\end{corollary}

\begin{proof}
Theorem~\ref{thm:stepseries}(3) with $z=s$ real, where $L_w(s)$ is the real window error.
\end{proof}

For a step length $s$, the largest window $p$ for which such a certificate exists is the largest window whose own loss one step of that length is certified to reduce; it is a limit of certification, not a proof that the step fails for longer windows, and the set of certified windows need not be an interval.
Since a window's own loss always decreases for a short enough step, this is a statement about step length and nonlinearity; the same series with $W_v$ in place of $W$ gives $V_v(z)$ and certifies a decrease of the error at chosen horizons.
$R_p$ is a lower bound on the radius of convergence; beyond the guaranteed disc $|z|<R_p$ the theorem says nothing.

\paragraph{What limits a single step: the radius, not the linearity.}
Inside the disc of convergence, the step series describes the effect of the step exactly to any chosen order, with a remainder that decreases like $(|z|/R)^{M+1}$; that the error is not linear in the step is no obstacle there, only a matter of how many orders are computed, which grows as $|z|$ approaches the radius.
Beyond the radius of convergence, no expansion about the current parameters converges, at any order: the rollout, as a function of the complex step, has a singularity or grows without bound within that distance.
A single step along $\mathbf d$ therefore meets two limits of different kind: how much of the window error any step of length $s$ can remove at all, which is the roughness theory's single-step value certified by its Krylov values; and the radius of the expansion in the step, which bounds the step lengths for which any expansion about the current parameters describes the step.
Beyond the radius, only expansions about other points along the path (that is, several shorter steps) or a direct evaluation of the step, without a series or a certificate, remain.
By Lemma~\ref{lem:stepholo}, $R_p\ge r_u/(L_\theta\phi_p)$: the guaranteed radius in the step is the radius $r_u$ in state space around the predicted states divided by how far a unit step moves the rollout, which grows with the window through $\phi_p$. The expansion in the step stops being guaranteed when a step of that length could move some predicted state to the edge of the state-space neighbourhood on which the learned map is holomorphic; this ties the radius in parameter space to the radius in state space that limits the expansion in the injection errors (Section~\ref{sec:analytic}).

\paragraph{How the radius and the coefficients depend on the window.}
For an Euler model, $\Gt=\uu+\dt\,\Hop$, so on $\mathcal D_q$ one can take $L_u=1+\dt\,\bar\ell$ and $L_\theta=\dt\,\bar L_\theta$ with $\bar\ell>0$, $\bar\ell\ge\sup\|\nabla_{\bm\xi}\Hop\|$ and $\bar L_\theta\ge\sup|\partial_z\Hop|$ over the complex neighbourhoods.
Then $\phi_p=\big((1+\dt\bar\ell)^p-1\big)/(\dt\bar\ell)\le(e^{\bar\ell\tau}-1)/(\dt\bar\ell)$ with $\tau=p\dt$, so with $A=\bar L_\theta(e^{\bar\ell\tau}-1)/\bar\ell$
\[
L_\theta\phi_p\le A,\qquad R_p\ge\min\big(r_\theta,\ r_u/A\big).
\]
For given constants this bound does not depend on $\dt$, and the guaranteed radius of the step series shrinks at most exponentially in the physical length of the window; when $r_u/A\le r_\theta$, Theorem~\ref{thm:stepseries}(1) with $R$ near $r_u/A$ bounds the coefficient of order $m$ by $A^mr_u^{1-m}$.
The constants $\bar\ell$, $\bar L_\theta$, $r_u$, $r_\theta$ are suprema over complex neighbourhoods of a rollout of $p$ steps of the network trained at that $\dt$; they can depend on $\dt$, on the direction and on the window ($\bar\ell$ bounds the complex Jacobian, not a Lyapunov exponent). These are upper bounds; how the actual radius shrinks is a measurement.
Along a fixed direction the coefficients $\mathbf Y^{(m)}$ are directional derivatives in $z$, computable by Taylor-mode (jet) forward differentiation through the rollout without forming derivative tensors.

\section{The two radii}\label{sec:tworadii}

The expansion in the injection errors (Section~\ref{sec:analytic}) is limited by a radius in state space, and the expansion in the step (Section~\ref{sec:stepseries}) by a radius in parameter space.
For networks with SiLU activation both are set by the same objects, the poles of the activation; this section bounds the radius of convergence of the step series from below and from above in terms of them.

\paragraph{Where the learned map stops being holomorphic.}
$\sigma(a)=a/(1+e^{-a})$ is holomorphic on $\mathbb C\setminus\mathcal P$, $\mathcal P=\{p_m=i\pi(2m+1):m\in\mathbb Z\}$, with a simple pole at each $p_m$ of residue $p_m$ (near $p_m$, $1+e^{-a}=1-e^{-(a-p_m)}$).
The network $\Nn(\mathbf x;\theta)$ is an MLP computed through hidden pre-activations $a^{(l)}_\nu=W^{(l)}_\nu\mathbf h^{(l-1)}+b^{(l)}_\nu$, $\mathbf h^{(l)}=\sigma(\mathbf a^{(l)})$, layer by layer, with a linear output layer $W^{\rm out}\mathbf h+b^{\rm out}$, and the learned map is $\Gt(\uu)=\alpha\uu+\beta+\gamma\odot\Nn((\uu-\mu)/\mathrm{sd})$ (Euler: $\alpha=1$, $\gamma=\dt\,\mathrm{sf}$; direct: $\alpha=0$), with $\alpha,\beta,\gamma,\mu,\mathrm{sd}$ fixed (not trained).
Each hidden pre-activation is a holomorphic function of the complex state and parameters as long as all hidden pre-activations computed before it avoid $\mathcal P$.
For a real pre-activation $a$, the distance to $\mathcal P$ is $|a\mp i\pi|=\sqrt{a^2+\pi^2}\ge\pi$.

Along a step $\theta+z\Pi^{1/2}\mathbf d$ (Section~\ref{sec:stepseries}), let $a_{q,l,\nu}(z)$ be the hidden pre-activation of neuron $\nu$ in layer $l$ when the rollout of a given start computes $\uu_{q+1}(z)$ from $\uu_q(z)$.
Because the start, the parameters and the direction are real, $a_{q,l,\nu}(\bar z)=\overline{a_{q,l,\nu}(z)}$ wherever it is defined (Schwarz reflection), and $\mathcal P=\bar{\mathcal P}$: pole crossings come in conjugate pairs.
Order the hidden pre-activations as the rollout computes them: by step $q$, then by layer $l$.

\begin{theorem}[Lower bound on the radius in the step]\label{thm:radlower}
Let $\rho>0$ and $D_\rho=\{|z|<\rho\}$. Suppose that, in the order of computation and for every start, each hidden pre-activation $a$ up to step $p-1$, which is then holomorphic on $D_\rho$, satisfies
\[
|a(z)-a(0)|<\sqrt{a(0)^2+\pi^2}\qquad\text{for }z\in D_\rho .
\]
Then no hidden pre-activation meets $\mathcal P$ on $D_\rho$, the rollout $\uu_q(z)$, $q\le p$, and the window error $L_w(z)$ are holomorphic on $D_\rho$, and the Taylor series at $0$ of $\uu_q(z)$, $q\le p$, and of $L_w(z)$ converge on $D_\rho$.
If, for each $a$, $|a(z)-a(0)-a'(0)z|\le K_a|z|^2$ on $D_\rho$, the hypothesis holds whenever $|a'(0)|\rho+K_a\rho^2\le\sqrt{a(0)^2+\pi^2}$ for all $a$.
\end{theorem}

\begin{proof}
Induction in the order of computation. The first hidden pre-activation at step $0$ is affine in $z$, since the start, $\mu$ and $\mathrm{sd}$ are fixed and the parameters are affine in $z$. If all earlier hidden pre-activations are holomorphic on $D_\rho$ and avoid $\mathcal P$, the activations and the state $\uu_q(z)$ are holomorphic (hence finite) there, so the next hidden pre-activation is holomorphic on $D_\rho$, and the hypothesis keeps it at distance below $\mathrm{dist}(a(0),\mathcal P)$ from $a(0)$, hence off $\mathcal P$. $L_w$ is a bilinear form in the states. A function holomorphic on $D_\rho$ has a Taylor series converging on $D_\rho$. The last statement is $|a(z)-a(0)|\le|a'(0)||z|+K_a|z|^2$.
\end{proof}

To first order in $z$ the hypothesis reads $|z|<\sqrt{a(0)^2+\pi^2}/|a'(0)|$ for every hidden pre-activation, so the guaranteed radius is about
\begin{equation}\label{eq:radest}
\min_{\text{starts},\,q<p,\,l,\,\nu}\ \frac{\sqrt{a_{q,l,\nu}(0)^2+\pi^2}}{|a'_{q,l,\nu}(0)|},
\end{equation}
over hidden layers, with neurons with $a'(0)=0$ contributing $+\infty$; this is an estimate, neither an upper nor a lower bound, and Theorem~\ref{thm:radlower} is the exact statement.

\begin{theorem}[Upper bound on the radius in the step]\label{thm:radupper}
Fix a start, a step $q<p$, a neuron $\nu$ of the last hidden layer $l=l_{\rm last}$, a pole $p_m$, and write $a(z)=a_{q,l_{\rm last},\nu}(z)$. Let $z_0\in\mathbb C$, $0<\delta<|z_0|$, $\rho_1=|z_0|+\delta$, $\Omega=D(z_0,\delta)\cup D(\bar z_0,\delta)$, and suppose that for some $\rho_1'>\rho_1$:
\begin{enumerate}
\item for this start, every hidden pre-activation computed before $a$ (at steps below $q$, and at step $q$ in layers below $l_{\rm last}$) and every other neuron of the last hidden layer at step $q$ is holomorphic on $D_{\rho_1'}$ and avoids $\mathcal P$ there, and $a$ avoids $\mathcal P$ on $D_{\rho_1'}\setminus\Omega$;
\item $|a(z)-a(z_0)-a'(z_0)(z-z_0)|+|a(z_0)-p_m|<|a'(z_0)|\,\delta$ on $|z-z_0|=\delta$;
\item $\gamma\odot W^{\rm out}_{\cdot\nu}(z_1)\ne0$, where $W^{\rm out}(z)$ is the output layer at the parameters $\theta+z\Pi^{1/2}\mathbf d$ and $z_1$ is the zero of $a-p_m$ given below.
\end{enumerate}
Then $a-p_m$ has exactly one zero $z_1$ in $D(z_0,\delta)$, it is simple, $\uu_{q+1}(z)$ has simple poles at $z_1$ and $\bar z_1$ with residue vectors $\mathbf c$ and $\bar{\mathbf c}$, $\mathbf c\ne0$, and the Taylor series of $\uu_{q+1}$ at $0$ has radius of convergence at most $|z_1|<\rho_1$.
If moreover (1) holds for every start, the window ends at $q+1$ with $w_{q+1}>0$, and $\mathbf c^\top\mathbf c\ne0$, the window error $L_w(z)$ has poles of order two at $z_1$ and $\bar z_1$, and its Taylor series at $0$ has radius of convergence at most $|z_1|$.
\end{theorem}

\begin{proof}
By (1), $a$ is holomorphic on $D_{\rho_1'}$, which contains the closed disc $\bar D(z_0,\delta)$. On $|z-z_0|=\delta$, $f=a-p_m$ and $g(z)=a'(z_0)(z-z_0)$ satisfy $|f-g|\le|a(z)-a(z_0)-a'(z_0)(z-z_0)|+|a(z_0)-p_m|<|g|$, so by Rouché's theorem $f$ has exactly one zero in $D(z_0,\delta)$, counted with multiplicity: a simple zero $z_1$, with $a'(z_1)\ne0$. By the reflection symmetry, $\bar z_1$ is a simple zero of $a-\bar p_m$.
Near $z_1$, $\sigma(a(z))=p_m/(a'(z_1)(z-z_1))+$ holomorphic, and $\uu_{q+1}=\alpha\uu_q+\beta+\gamma\odot(W^{\rm out}\mathbf h+b^{\rm out})$ with all other terms holomorphic near $z_1$ by (1); so $\uu_{q+1}$ has a simple pole at $z_1$ with residue $\mathbf c=\gamma\odot W^{\rm out}_{\cdot\nu}(z_1)\,p_m/a'(z_1)\ne0$ by (3), and by reflection a simple pole at $\bar z_1$ with residue $\bar{\mathbf c}$.
By (1), $\uu_{q+1}$ is holomorphic on $D_{\rho_1'}$ minus the zeros of $a-p_{m'}$, $p_{m'}\in\mathcal P$, which lie in $\Omega$ and form a relatively closed discrete set; its complement in $D_{\rho_1'}$ is connected and contains $0$ ($0\notin\Omega$ because $\delta<|z_0|$, and $a(0)\notin\mathcal P$ because it is real). If the Taylor series at $0$ converged on a disc containing $z_1$, its sum would be holomorphic near $z_1$ and, by the identity theorem on the connected set where both are defined, equal to $\uu_{q+1}$ there, which is unbounded near $z_1$: a contradiction.
For the window error, the terms of horizons $\le q$ and those of the other starts are holomorphic near $z_1$ by (1) for every start, and the horizon-$(q+1)$ term of this start is $\tfrac{w_{q+1}}{2n}\ee_{q+1}^\top\ee_{q+1}$ with $\ee_{q+1}=\uu_T(t_{q+1})-\uu_{q+1}$, which has a pole of order two at $z_1$ with coefficient $\tfrac{w_{q+1}}{2n}\mathbf c^\top\mathbf c\ne0$ (and at $\bar z_1$ with the conjugate coefficient); the same argument applies.
\end{proof}

For windows longer than $q+1$ the later steps evaluate the network at a state with a pole, which can produce a non-isolated singularity; nothing is claimed for them.
The claim about $\uu_{q+1}$ needs (1) only for the given start.

\begin{corollary}[Two-sided form near the linear prediction]\label{cor:radtwo}
In the setting of Theorem~\ref{thm:radupper}, assume $a'(0)\ne0$, let $z_0=(p_m-a(0))/a'(0)$, where the linear prediction $a(0)+a'(0)z$ meets $p_m$, assume $a$ is holomorphic on $\bar D(0,\tfrac32|z_0|)$, and let $K>0$ with $K\ge\sup_{|z|\le\frac32|z_0|}|a(z)-a(0)-a'(0)z|/|z|^2$.
If $K|z_0|\le|a'(0)|/8$ and hypotheses (1) and (3) hold with $\delta=4K|z_0|^2/|a'(0)|$, then the radius of convergence of the Taylor series of $\uu_{q+1}$ at $0$ is less than
\[
|z_0|+\frac{4K|z_0|^2}{|a'(0)|},\qquad |z_0|=\frac{|p_m-a(0)|}{|a'(0)|}\ \Big(=\frac{\sqrt{a(0)^2+\pi^2}}{|a'(0)|}\text{ for }p_m=\pm i\pi\Big).
\]
\end{corollary}

\begin{proof}
$a(z)-p_m=\big(a(z)-a(0)-a'(0)z\big)+a'(0)(z-z_0)$. With $g(z)=a'(0)(z-z_0)$, on $|z-z_0|=\delta\le|z_0|/2$ one has $|f-g|\le K(|z_0|+\delta)^2\le\tfrac94K|z_0|^2<4K|z_0|^2=|g|$, so Rouché's theorem gives exactly one zero of $a-p_m$ in $D(z_0,\delta)$; the rest is the proof of Theorem~\ref{thm:radupper} with this disc.
\end{proof}

\paragraph{The link between the radii.}
By the chain rule along the step, $a'_{q,l,\nu}(0)=\nabla_\uu a_{q,l,\nu}\cdot\mathbf Y^{(1)}_q+\partial_\theta a_{q,l,\nu}\,\Pi^{1/2}\mathbf d$, with $\mathbf Y^{(1)}_q=\Psih_q\mathbf d$ the rollout feature along the step, and by Cauchy--Schwarz the estimate \eqref{eq:radest} satisfies, exactly,
\[
\eqref{eq:radest}\ \ge\ \min_{\text{starts},\,q<p,\,l,\,\nu}\frac{\sqrt{a_{q,l,\nu}(0)^2+\pi^2}}{\|\nabla_\uu a_{q,l,\nu}\|\,|\Psih_q\mathbf d|+|\partial_\theta a_{q,l,\nu}\,\Pi^{1/2}\mathbf d|}.
\]
The quantity $r^{\rm lin}_{q}=\min_{l,\nu}\sqrt{a_{q,l,\nu}^2+\pi^2}/\|\nabla_\uu a_{q,l,\nu}\|$ is, to first order, the distance in complex state space from the predicted state to the states at which some neuron reaches a pole, that is, a first-order version of the state-space radius $r_u$ of Assumption~\ref{ass:stepanalytic} around the predicted states.
So, to first order and when the term through the state is the larger, the radius in parameter space is at least the radius in state space divided by how far a unit step moves the predicted state, $|\Psih_q\mathbf d|$.
The pre-activation pole distance $\sqrt{a^2+\pi^2}$ is at least $\pi$ independently of the window, and $r^{\rm lin}_q\ge\pi/\max_{l,\nu}\|\nabla_\uu a_{q,l,\nu}\|$ over the predicted states visited; the window enters through the growth of $|\Psih_q\mathbf d|$ with the propagators, through the gradients $\nabla_\uu a$ at the predicted states visited, and through the larger set of steps over which the minimum is taken.
The last three sentences are first-order readings of the theorems, not further theorems; the upper bound applies in addition only under the hypotheses of Theorem~\ref{thm:radupper}, which involve values of the pre-activations on complex discs and have not been verified for trained networks.

\section{The best number of physical steps for one training step}\label{sec:bestp}

\paragraph{The criterion.}
Windows of different lengths have different losses, so they are compared on a fixed target that does not depend on the window.
Let $\mathcal H$ be a set of held-out starts and $N$ a number of physical steps, and let the target be the held-out window error
\[
V(\theta)=\tfrac12\,\ee_{\mathcal H}^\top W_N\,\ee_{\mathcal H}=\frac1{N}\sum_{q=1}^N\frac1{|\mathcal H|}\sum_{k\in\mathcal H}\tfrac12|\ee_k(t_q)|^2,
\]
so that $\sqrt{2V}$ is the root-mean-square rollout error over all pairs of a held-out start and a horizon $q\le N$.
For a training window of $p$ physical steps on the training starts $\mathcal T$, with weights $w_q\ge0$, $q\le p$, and the average $1/|\mathcal T|$ over training starts in $W_p$, let $g_p=-\Psih_{\mathcal T}^\top W_p\ee_{\mathcal T}$ be the rescaled gradient of its loss $L_p$; assume $g_p\ne0$, let $\mathbf d_p=-g_p/|g_p|$ and, for a fixed rescaled step length $s>0$, $\theta_p(s)=\theta+s\,\Pi^{1/2}\mathbf d_p$.
For a set $P$ of at least two candidate windows,
\[
p^\ast(s)\in\arg\min_{p\in P}V\big(\theta_p(s)\big)
\]
is the best number of physical steps for this training step.
At a fixed rescaled step length the window enters only through the direction $\mathbf d_p$: choosing $p$ is choosing one of finitely many directions of length $s$.
With a fixed step size $\alpha$ instead, as in the expected criterion of Section~\ref{sec:operators}, the length $\alpha|g_p|$ also depends on the window, and the two criteria can rank windows differently.

\begin{proposition}[Alignment with the held-out target]\label{prop:align}
With $g_V=\Pi^{1/2}\nabla_\theta V=-\Psih_{\mathcal H}^\top W_N\ee_{\mathcal H}$, exactly
\[
\frac{d}{ds}V\big(\theta_p(s)\big)\Big|_{s=0}=\langle g_V,\mathbf d_p\rangle=-\frac{\langle g_V,g_p\rangle}{|g_p|},\qquad
\langle g_V,g_p\rangle=\big(W_N\ee_{\mathcal H}\big)^\top\Psih_{\mathcal H}\Psih_{\mathcal T}^\top\big(W_p\ee_{\mathcal T}\big),
\]
the inner product of the weighted held-out and training errors through the cross rollout kernel between held-out and training starts; for $g_V\ne0$ the derivative is $-|g_V|\cos(g_V,g_p)$.
\end{proposition}

\begin{proof}
Chain rule with $\theta_p'(0)=\Pi^{1/2}\mathbf d_p$, and the exact form of the gradient of a rollout loss (Section~\ref{sec:grad}) applied separately to the held-out and the training rollouts.
\end{proof}

\begin{proposition}[Series, certificate and ranking]\label{prop:bestp}
For each $p\in P$ let Assumption~\ref{ass:stepanalytic} hold for the direction $\mathbf d_p$ on the held-out rollouts of $N$ steps, with constants $r_u,r_\theta,L_u$ common to all $p$ and $L_{\theta,p}$ depending on $p$ (possible for a finite $P$ by taking the smallest radii and the largest $L_u$), guaranteed radius $R^V_p=\min(r_\theta,r_u/(L_{\theta,p}\phi_N))$, and let $\nu_m(p)$ be the Taylor coefficients of $z\mapsto V(\theta+z\Pi^{1/2}\mathbf d_p)$, given by Theorem~\ref{thm:stepseries}(2) with $W_N$, the held-out errors and the sensitivities along $\mathbf d_p$.
\begin{enumerate}
\item For $s<\rho<R^V_p$ and every $M\ge1$, $V(\theta_p(s))$ lies in the interval $\sum_{m=0}^M\nu_m(p)s^m\pm B^V_p(\rho)\frac{(s/\rho)^{M+1}}{1-s/\rho}$, with $B^V_p(\rho)=\|\ee_{\mathcal H}\|_{W_N}L_{\theta,p}\rho\|\phi\|_{W_N}+\tfrac12L_{\theta,p}^2\rho^2\|\phi\|_{W_N}^2$; $p$ is certified best in $P$ when the upper end of its interval lies below the lower ends of all others.
\item Let $a_p=\langle g_V,g_p\rangle/|g_p|$, so $\nu_1(p)=-a_p$, let $p_1$ maximize $a_p$ over $P$ and $\delta=a_{p_1}-\max_{p\ne p_1}a_p$. For $\rho_p\in(0,R^V_p)$ let $c_p=2B^V_p(\rho_p)/\rho_p^2$. If $0<s\le\min_p\rho_p/2$ and $s\,\max_pc_p<\delta/2$, then $\arg\min_{p\in P}V(\theta_p(s))=\{p_1\}$.
\end{enumerate}
\end{proposition}

\begin{proof}
(1) is Theorem~\ref{thm:stepseries}(3) for $V$ along each $\mathbf d_p$ (the theorem holds for any unit direction, starts and weights); the radius involves $\phi_N$, the evaluation length, because the target is a rollout of $N$ steps.
(2) With $M=1$ and $s\le\rho_p/2$, $|V(\theta_p(s))-V+s\,a_p|\le B^V_p(\rho_p)(s/\rho_p)^2/(1-s/\rho_p)\le c_ps^2$. Hence $V(\theta_{p_1}(s))\le V-s\,a_{p_1}+c_{p_1}s^2$ and, for $p\ne p_1$, $V(\theta_p(s))\ge V-s(a_{p_1}-\delta)-c_ps^2$; the first is strictly smaller when $s\,(c_{p_1}+c_p)<\delta$, which the hypothesis gives. The hypothesis forces $\delta>0$, so $p_1$ is the unique maximizer.
\end{proof}

\paragraph{What decides the best window.}
By Proposition~\ref{prop:align} the first-order part of the decrease is the alignment $a_p$ of the training window's step direction with the held-out target gradient, through the cross rollout kernel; by Proposition~\ref{prop:bestp} the higher orders along each direction enter the ranking, and decide it when the first-order gap $\delta$ is too small for the step length.
We expect, but have not shown, that longer training windows can align better with a target averaged over many steps and that their directions excite more strongly curved directions of the held-out rollout; both can be measured through $a_p$ and the coefficients $\nu_m(p)$.
No general monotonicity in $p$ follows.
The guaranteed radius $R^V_p$ is set by the evaluation length $N$ through $\phi_N$; the window enters only through $\mathbf d_p$, via all constants of Assumption~\ref{ass:stepanalytic} (here only $L_{\theta,p}$ is allowed to vary), and $R^V_p$ is only a lower bound on the radius of convergence, whose actual value depends on the direction.

\paragraph{The alignment as a series in the injection errors.}
The rate series of Proposition~\ref{prop:rateorders} and Theorem~\ref{thm:tradeoffbound} were stated for one set of starts.
They apply here with $\mathcal H\cup\mathcal T$ as one set of starts and start-dependent weights: $W_v=W_N$ on the held-out block and $0$ on the training block, $W=W_p$ on the training block and $0$ on the held-out block, rollout length $\max(N,p)$, and the injections of both sets scaled by the same $s$; Assumption~\ref{ass:analytic} and the admissibility of some $\rho>1$ must then hold for every start in $\mathcal H\cup\mathcal T$, out to the longest candidate window, which can fail at long windows.
The bound $B(\rho)$ keeps its product form, the two factors being averaged over $\mathcal H$ and $\mathcal T$ separately.
Ranking by $a_p=\langle g_V,g_p\rangle/|g_p|$ also needs $|g_p|^2=\langle g_p,g_p\rangle$, which is a series of the same kind with $v=w$ on the training starts, and whose square root is analytic while $|g_p|$ stays away from zero.
Every coefficient of these series lies along the true trajectories and is nested in $p$, so one true rollout of the longest window serves every candidate; the leading term $\langle\nabla^{(1)}_V,\nabla^{(1)}_p\rangle$ still involves the first-order gradient of each candidate.
The best window is a property of the current parameters; nothing is claimed about later training steps, and a window is shown worse than another only by such a comparison.

\section{The curvature of the multistep loss}\label{sec:curvature}

The Hessian of $L_w$ in rescaled coordinates is exactly $\Psih^\top W\Psih-\sum_{k,q}(W\ee)_{k,q}\cdot\Pi^{1/2}\partial_\theta^2\uu_{p,k}(t_q)\Pi^{1/2}$.
Its Gauss--Newton part $\Psih^\top W\Psih$ has the nonzero eigenvalues of $\Psih\Psih^\top W$, the rollout kernel of the rollout-kernel note.
That note proves an upper bound on the largest eigenvalue of this Gauss--Newton part (its rate theorem): for Euler models, at horizon $\tau_q$ the endpoint kernel has largest eigenvalue at most $e^{2|\mu|\dt}\varphi_\mu(\tau_q)^2\max_is_i^2$, with $\varphi_\mu(\tau)=(e^{\mu\tau}-1)/\mu$, $s_i^2$ the largest one-step rates of $\Hop$, and $\mu$ a bound on $\mu_{\log}(\nabla\Hop)+\tfrac{\dt}{2}\|\nabla\Hop\|^2$ at every visited predicted state ($\mu_{\log}$ the logarithmic norm); for the window loss the bound is the $w_q$-weighted sum of these.
It also proves that, at fixed parameters and start and under existence and smoothness of the learned ODE and its sensitivity equation, the rollout kernel of an Euler model converges as $\dt\to0$ at fixed horizon; for direct models only the upper bound, which grows with the number of steps, is available.
The growth of the Gauss--Newton curvature with the window is therefore a property of the learned network at the training starts; its exponential growth on Lorenz-63 and the $\dt^{-2}$ growth for direct models are measured (Section~\ref{sec:numerics}).
Near a minimum, linearized gradient descent is stable only for steps below $2$ over the largest Hessian eigenvalue; choosing the step length from the local quadratic model along the update direction is a heuristic.

\section{Lorenz-63 measurements}\label{sec:numerics}

These are measurements on one system (Euler and direct models, $\dt\le0.08$, MLP $3\times256$ SiLU unless stated, 3 seeds); they illustrate but do not support the proofs above.
\begin{itemize}
\item The exact identities (features lemma, error recursion and its orders, the exact effect of a parameter step) hold to $10^{-11}$--$10^{-12}$.
\item The Gauss--Newton curvature at the start of fine-tuning (for separately trained models, so the fixed-parameter limit theorem does not predict this) is the same function of the window for all $\dt$ for Euler models ($\approx2.6\times10^3$, $10^5$, $4\times10^7$, $6.5\times10^9$ for windows $0.32,0.64,1.28,2.56$) and grows like $\dt^{-2}$ for direct models. With a fixed learning rate windows $\ge1.28$ fail to train; with a step size matched to the curvature windows up to $2.56$ train.
\item The size of the gradient's higher-order part, $|g_w-\nabla^{(1)}_w|/|\nabla^{(1)}_w|$ (Euler, $\dt=0.08$, widths 64--512, 1024 held-out starts), is about $10^{-4}$ for windows $\le0.32$, $10^{-3}$ at $0.64$ and $0.01$--$0.18$ between $0.96$ and $1.44$; it jumps between $1.52$ and $1.60$ in every model.
Removing the $1\%$ of starts with the largest first-order error (10 starts) brings it below $0.2$, and removing $5\%$ below $0.05$, up to $2.48$: a few percent of the starts carry most of the higher-order part of the gradient beyond about $1.5$.
One to five starts leave the estimated (not certified, and taken at held-out rather than the rollout's true states) radius of Remark~\ref{rem:radius} at windows coinciding with the jump (one seed per model).
\item During plain gradient descent from trained models, the sign of the exact rate of Lemma~\ref{lem:rates}, evaluated on held-out starts, agrees with the sign of the measured change of the error at horizons up to $1.92$ in about $80$--$98\%$ of ten-step intervals; the accumulated size is not predicted (the steps are not small).
\item Measured by the root-mean-square over starts, windows $1.28$--$2.56$ give the lowest errors at their own horizons; measured by the median or 90th percentile, short windows ($0.32$--$0.64$) are best at every horizon up to $2.56$. The root-mean-square is about $3$--$5$ times the median at horizon $0.64$ and about $20$ times at $1.28$, at every width from 64 to 512, consistent with a heavy tail of the error over starts.
\end{itemize}

\section{Status and open questions}\label{sec:status}

\begin{center}
\small
\begin{tabular}{p{0.45\textwidth}p{0.47\textwidth}}
\toprule
Statement & Status \\
\midrule
Error recursion (Lemma~\ref{lem:recursion}), exact rates (Lemma~\ref{lem:rates}) & exact; the rates under gradient flow with fixed $\Pi$, or with the actual update direction \\
Orders of the error, features and gradient (Definition~\ref{def:orders}, \eqref{eq:Phij}--\eqref{eq:gradorders}), every order of the rate (Proposition~\ref{prop:rateorders}) & explicit formal series \\
Convergence of the error series (Theorem~\ref{thm:convergence}) & proved under Assumption~\ref{ass:analytic} and the hypothesis of Theorem~\ref{thm:allorders}; uses Banach's theorem on symmetric forms \\
Convergence and remainder of the rate, geometric decay of the orders (Theorem~\ref{thm:tradeoffbound}) & proved under Assumption~\ref{ass:analytic} for admissible $\rho>1$; the $N=2$ criterion needs $\rho>2$ \\
Bounds on $\ee-\ee^{(1)}$ and $g_w-\nabla^{(1)}_w$, certified sign (Theorems~\ref{thm:allorders}, \ref{thm:gradbound}, Corollary~\ref{cor:sign}) & proved under Assumption~\ref{ass:analytic} and the hypothesis of Theorem~\ref{thm:allorders}; checkable only with certified constants \\
Window of controlled nonlinearity $\tau_{\rm nl}$ (Corollary~\ref{cor:closed}) & proved under uniform constants of the trained network; controls the size of the higher orders, not descent \\
Floor during multistep training, rollout floor (Proposition~\ref{prop:floorms}, Lemma~\ref{lem:inverse}, Theorem~\ref{thm:rolloutfloor}) & proved without expansion under the roughness paper's hypotheses, inputs the window's true states, $\ell$ a~posteriori \\
Floor as the supremum over all seminorms, certified by Krylov values (\eqref{eq:supfloor}, Theorem~\ref{thm:rolloutfloor}) & proved (roughness paper's supremum and Krylov theorems); fixed seminorms give weaker bounds \\
Direct models, unbuilt identity (Corollary~\ref{cor:directrollout}) & proved; exact for the step taken with step-averaged features, first order in the step with current features; uses $\|S_h-\mathrm{id}\|\le M_1\dt$ \\
Kernel-level results for rollouts & rates and the every-step bound exact; with step-averaged features an exact bound on the step taken; best-step value, slow part and Krylov otherwise first order in the step, no remainder bound \\
Upper bound on the Gauss--Newton part of the curvature & proved in the rollout-kernel note; exponential growth and the direct $\dt^{-2}$ growth measured \\
How far a window can usefully be extended & not determined: certificates only, no converse \\
Concentration of the gradient on a few starts & measured \\
Floor at a fixed horizon from the roughness of $S_\tau$ & open \\
Rollout error as true separation plus model error, learned versus true propagators along the true trajectory (Lemma~\ref{lem:split}, Proposition~\ref{prop:propsplit}) & exact; bounds under true-system constants $C_S,\Lambda_S$ \\
Roughness theory for each operator error, cross-kernels (Theorem~\ref{thm:operators}) & proved, with the Frobenius inner product on tensor-valued sampled functions; same exactness as for the values (best-step value first order in the step with current features) \\
Local complexity of operator targets (Proposition~\ref{prop:opcomplexity}) & proved under the roughness paper's Assumption~1 and smoothness of the flow, with $\bar M^{(k)}$ in the Frobenius norm \\
Norms do not determine the sign (Proposition~\ref{prop:normsign}) & proved for hypothetical injections with coefficients held fixed, total norm only, small injections, $v\ne w$ \\
Expected rate and the window for one step (Proposition~\ref{prop:expectedrate}) & \eqref{eq:onestepdecrease} exact; series under almost-sure admissibility for a fixed $\rho>1$ and $\mathbb E[B(\rho)]<\infty$; the moments are an input \\
Sensitivities of every order along a step (Lemma~\ref{lem:sens}) & exact \\
Step series, coefficient bounds, remainder, certified decrease, largest certified window (Lemma~\ref{lem:stepholo}, Theorem~\ref{thm:stepseries}, Corollary~\ref{cor:stepcap}) & proved under Assumption~\ref{ass:stepanalytic}; constants of the trained network, not certified yet; a limit of certification, no converse \\
Lower and upper bounds on the radius in the step from the poles of the activation (Theorems~\ref{thm:radlower}, \ref{thm:radupper}, Corollary~\ref{cor:radtwo}) & proved for SiLU MLPs; the upper bound under hypotheses (1)--(3) (values on complex discs, not verified for trained networks), for last-layer neurons, for the state one step later and for the window ending there ($\mathbf c^\top\mathbf c\ne0$); the radius as at least the state radius over the state displacement per unit step is a first-order reading (a lower bound on an estimate) \\
Lower bound on the step-series radius (Euler) & proved; independent of $\dt$ for given network constants; how the actual radius shrinks is to be measured \\
\bottomrule
\end{tabular}
\end{center}

\paragraph{Open questions.}
What training reaches over a run; the effect of the random sampling matrix of minibatches; adaptive optimizers beyond the instantaneous rate; a description of starts whose error exceeds the radius; and a local-complexity floor at a fixed horizon.
Evaluating $r$, $D_i$, $D_i'$, $F_i$ for trained networks would make Corollary~\ref{cor:sign} a checkable certificate.

\end{document}
